% =============================================================================
%  Paper A — IEEE Transactions on Mobile Computing (TMC) submission template
%
%  Title: WiFi CSI Cross-Domain Gesture Recognition:
%         A Reproducibility Audit and the Boundary of Multi-Rx Representation
%
%  Generated from: docs/叙事A_章节草稿_v2.md  (Version 4, 2026-09-07)
%  Companion paper draft: docs/叙事A_论文大纲.md
%  Evidence base:        docs/findings_B1-B14.md (B1-B20)
%  Related Work source:  docs/literature_review.md (v2)
%  Figures source:       scripts/make_paper_figs.py + scripts/make_fig0_pipeline.py
%                        + scripts/make_fig5_mechanism.py
%
%  Compile with: pdflatex paper_ieee_tmc.tex  (twice for cross-references)
%                bibtex paper_ieee_tmc  (after first pdflatex pass)
%                pdflatex paper_ieee_tmc.tex  (twice more)
% =============================================================================
\documentclass[11pt]{article}

% Originally written for IEEE TMC submission.
% Converted to arXiv-compatible article class.
% IEEEtran-specific commands and constraints removed.

% Stub IEEE keywords environment for arXiv article class compatibility.
\newenvironment{IEEEkeywords}
  {\par\noindent\textbf{Index Terms:}\ }
  {\par}

% ----------  packages  ----------
\usepackage[utf8]{inputenc}
\usepackage[T1]{fontenc}
\usepackage{amsmath,amssymb,amsfonts}
\usepackage{amsthm}        % theorem environments (added for §3.5 Statistical Foundation)
\usepackage{algorithmic}
\usepackage{algorithm}
\usepackage{graphicx}
\usepackage{textcomp}
\usepackage{xcolor}

% Theorem environments for §3.5 Statistical Foundation.
\newtheorem{theorem}{Theorem}
\newtheorem{lemma}[theorem]{Lemma}
\newtheorem{corollary}[theorem]{Corollary}
\theoremstyle{definition}
\newtheorem{definition}{Definition}

\usepackage{booktabs}
\usepackage{multirow}
\usepackage{array}
\usepackage{url}
\usepackage{cite}
\usepackage{hyperref}
\usepackage{bm}
\usepackage{siunitx}
\usepackage{subcaption}
\usepackage{microtype}   % better justification / hyphenation

% ----------  hyperref settings  ----------
\setlength{\emergencystretch}{3em}   % avoid overfull lines
\hypersetup{
    colorlinks=true,
    linkcolor=blue,
    citecolor=blue,
    urlcolor=blue,
    pdfauthor={Anonymous},
    pdftitle={WiFi CSI Cross-Domain Gesture Recognition: A Reproducibility Audit Revealing Dataset-Conditioned Optimality of the Multi-Rx Representation},
}

% ----------  custom commands  ----------
\newcommand{\BVP}{\text{BVP}}
\newcommand{\CSI}{\text{CSI}}
\newcommand{\MDE}{\text{MDE}}
\newcommand{\LODO}{\text{LODO}}
\newcommand{\Rx}{\text{Rx}}
\newcommand{\MultiRx}{Multi-Rx}
\newcommand{\CI}{\text{CI}}
\newcommand{\PR}{\text{PR}}
\newcommand{\pp}{\text{pp}}
\newcommand{\ie}{\textit{i.e.}}
\newcommand{\eg}{\textit{e.g.}}
\newcommand{\etal}{\textit{et al.}}
\newcommand{\TODO}[1]{[TODO removed: #1]}

% ----------  title  ----------
\title{WiFi CSI Cross-Domain Gesture Recognition:\\
       A Reproducibility Audit Revealing Dataset-Conditioned Optimality\\
       of the Multi-\Rx{} Representation}

\author{%
    Anonymous Authors\\[2pt]
    \small\textit{Anonymous Institution}
}

\markboth{Anonymous Preprint}%
{Anonymous et al.}

\begin{document}

\maketitle

% =============================================================================
\begin{abstract}
WiFi CSI-based cross-domain gesture recognition has attracted intensive research in the past six years, yet reported improvements rarely reproduce under controlled evaluation. We present a \textbf{reproducibility audit} covering twelve paired comparisons across ten method axes---mixup augmentation, loss reweighting, domain adversarial training, snapshot ensembling, attention variants, multi-scale hop settings, backbone-capacity controls, and BVP-representation alternatives (multi-\Rx, RMS-aggregation, single-\Rx)---on the Widar3.0 benchmark (6 leave-one-domain-out folds; primary comparisons use $n = 18$--$30$ strict paired samples depending on the overlap of seed sets with the baseline; the C10 larger-backbone capacity control uses $n = 18$ strict paired samples from 6 folds $\times$ 3 shared seeds and is additionally evaluated on a 30-run set across 5 seeds as informational; 4 folds $\times$ 5 seeds = 20 strict paired samples on MMFi). Four of the audited algorithmic comparisons \textbf{fail to reach significance} in our primary paired evaluations ($p > 0.45$, $|d| < 0.3$). We provide closed-form statistical foundations (Section~\ref{sec:statfoundation}) for the audit and show that single-seed reports of $\le 2$\,pp gain are \emph{provably underpowered} under any paired design at our sample budget ($n \ge 56$ required, $n \le 18$ available). Variance decomposition reveals the root cause: within-seed standard deviation can exceed 4\,pp, which is comparable to the largest reported effect and makes single-seed studies unable to distinguish real gains from noise. We then introduce a minimal signal-representation modification---expanding the BVP operator to retain the multi-Rx antenna dimension (comparison \textbf{C6})---that yields \textbf{+8.14\,pp} (Cohen's $d = 1.71$, all 18 paired comparisons positive). We further test this finding on a second dataset MMFi (27 gestures, 4 environments; comparison \textbf{C7}, v1 path-labeled cache, 1080 samples) and find that the multi-\Rx\ advantage does \textbf{not transfer}: multi-scale is nominally ahead on MMFi ($+0.93$\,pp vs.\ multi-\Rx, $p = 0.116$; $+1.84$\,pp vs.\ RMS-agg, $p = 9.2\times10^{-4}$). t-SNE visualization reveals the geometric root cause: MMFi's four environments are indistinguishable in the BVP feature space, while Widar3.0 retains partial gestural clustering. \textbf{No single BVP representation is universally optimal; the choice must be matched to the dataset's geometric structure.} Mechanistic analysis (Section~\ref{sec:mechanism}) attributes the gain to an effective increase in the antenna participation ratio (1.0 $\to$ 1.89) and a $1.82\times$ (${\sim}82\%$) increase in total Fisher information, both consequences of front-end signal preservation rather than back-end algorithmic novelty. \emph{As an additional honesty check}, we re-evaluated the larger-backbone capacity control (C10, CBAM+ResNet18 vs LeNetCSI\_Attn) on a 50{,}000-sample subset of the dedup-ed Widar3.0 pool (30.55\,\% of the full 163{,}650 unique link-level samples) under the identical paired LODO protocol: the \emph{direction is stable} across subsets---the smaller backbone wins on both (12k: $\Delta = -4.61$\,pp, $p = 0.028$, 3/18 wins for CBAM+ResNet18; 50k: $\Delta = -5.75$\,pp, $p = 0.0609$, 4/18 wins, Cohen's $d = -0.44$, 95\,\% CI $[-11.77, +0.26]$)---but the \emph{significance does not survive}: the paired-difference standard deviation inflates from 8.9 to 13.0\,pp and the 50k 95\,\% CI crosses zero with $p$ marginally above $\alpha = 0.05$. We therefore report the backbone comparison as direction-consistent yet statistically fragile across subsets---a variance-inflation effect that motivates the 30\,\%-subset policy in Section~\ref{sec:boundary}---not as evidence of a winner swap. Our findings reframe cross-domain CSI research from a model-architecture problem to a joint (representation $\times$ protocol $\times$ dataset-difficulty) problem, and provide a paired evaluation protocol (Section~\ref{sec:methodology}) that other researchers can adopt to make their own reported gains auditable.
\end{abstract}

\begin{IEEEkeywords}
WiFi sensing, channel state information (CSI), cross-domain gesture recognition, reproducibility, paired multi-seed evaluation, minimum detectable effect (MDE), antenna diversity.
\end{IEEEkeywords}

% =============================================================================
\section{Introduction}
\label{sec:intro}

WiFi Channel State Information (\CSI) cross-domain gesture recognition has emerged as a promising paradigm for contactless human-computer interaction, with applications spanning smart homes, elderly care, and human-computer interfaces~\cite{zheng2019widar, wang2022csi}. The task is challenging because the same gesture performed by the same user yields substantially different \CSI{} signatures under varying environments, user poses, and orientations. Since 2019, the community has proposed dozens of methods---domain adversarial training~\cite{sheng2020}, attention-based architectures~\cite{zhang2021}, self-supervised pretraining~\cite{ma2022}, multimodal fusion~\cite{xie2023}---and the field has grown rapidly.

Yet a disturbing pattern has emerged: \textbf{reported gains rarely reproduce}. To our knowledge, in our preliminary survey of 47 \CSI{} cross-domain papers published in top venues between 2020 and 2025, only 3 reported variance across multiple random seeds; few reported effect-size confidence intervals or cross-validated their methods on a second dataset. This is reminiscent of the broader reproducibility crisis in machine learning~\cite{sculley2018, pineau2021}, but with a key difference: \CSI{} cross-domain researchers routinely publish ``significant'' results from single-seed experiments without questioning their robustness.

We argue that this crisis is not an accident. It is the predictable consequence of an effect-size landscape in which the \textbf{training-seed standard deviation is larger than the typical reported improvement}~\cite{henderson2018, pineau2021, dror2018}. In such a regime, any single-seed experiment is essentially a random draw from a noisy distribution, and the chance of falsely claiming a ``significant'' gain is high.

Before substantiating this claim, we cross-validate it against independently published work to ensure we are not measuring an artifact of our protocol:
\begin{itemize}
    \item \textbf{Reality Check \#1 (SenseFi benchmark, 2023):} Yang \etal\ published a benchmark finding that ``shallow models outperform deep models'' on WiFi sensing tasks~\cite{yang2023sensefi}. We did not re-run their exact shallow-vs-deep pairing; instead we re-tested the nearest architectural claim in our audit---\textbf{C1}, MHA vs.\ convolutional attention on Widar3.0 ($n = 18$ strict paired samples): the apparent $+1.13$\,pp MHA advantage is not statistically distinguishable from noise under paired evaluation ($p = 0.478$, $d = 0.17$), consistent with SenseFi's broader observation that architectural sophistication does not reliably translate into cross-domain gains.
    \item \textbf{Reality Check \#2 (CSI-Bench, NeurIPS 2025):} Zhu \etal\ released a 461-hour, 35-user, 16-device benchmark and reported all results as \emph{3-seed mean\,$\pm$\,std}~\cite{zhu2025csibench}. Our MDE analysis shows that 3-seed designs have $\sim$70\% power to detect 3\,pp effects but \textless 25\% power for 1\,pp effects---the very effect range where most reported gains live.
    \item \textbf{Reality Check \#3 (RuView open-source deployment, 2026):} Cohen's ESP32-mesh-based WiFi sensing stack~\cite{cohen2026ruview} is among the most engineered industrial systems publicly available; per its public documentation, its maintainers report a substantial drop in accuracy when deployed across environments (``100\% presence detection retracted to 82.3\% under distribution shift'', per the project's README)---a real-world confirmation that \CSI{} features are \emph{distribution-locked}.
\end{itemize}

These three independent checks---peer-reviewed benchmark, large-scale NeurIPS benchmark, and an open-source industrial system---converge on the same conclusion: \CSI{} cross-domain evaluation as currently practiced yields noisy, non-reproducible claims, even when researchers use industry-standard protocols. Our contribution is therefore not to \emph{describe} this crisis but to \emph{demonstrate} it under a 5-seed paired protocol with pre-registered statistical power, and to identify a representation intervention that survives the audit while remaining bounded by dataset separability.

To verify the hypothesis rigorously, we constructed a \textbf{paired cross-domain benchmark} (Section~\ref{sec:methodology}) on the standard Widar3.0 dataset~\cite{zheng2019widar} and systematically re-evaluated twelve paired comparisons across ten method axes (Section~\ref{sec:results}). Several of the audited methods \textbf{fail to reach statistical significance} under paired tests. Variance decomposition (Section~\ref{sec:variance}) shows that within-seed standard deviation can reach several percentage points, comparable to typical reported effects. Even more critically, a Minimum Detectable Effect (\MDE) analysis (Section~\ref{sec:mde}) makes the power of every comparison computable in advance: at our sample budget the paired design can only detect effects above roughly 1--5\,pp (depending on within-pair variance), so the observed failures must be read as ``not distinguishable from noise at our measurement precision,'' not as proofs of absence.

Where, then, can the field move forward? We identify a minimal signal-representation modification---\textbf{expanding the \BVP{} operator to retain the multi-\Rx{} antenna dimension}---that achieves \textbf{+8.14\,pp} with Cohen's $d = 1.71$ across all 18 paired comparisons of \textbf{C6} (Section~\ref{sec:mechanism}). Mechanistic analysis (Section~\ref{sec:mechanism}) attributes the gain to two complementary effects: an effective increase in the antenna participation ratio from $1.00$ to $1.89$ (Section~\ref{sec:mechanism}), and a $1.82\times$ (${\sim}82\%$) increase in total Fisher information. Both are consequences of front-end signal preservation rather than back-end algorithmic novelty.

However, we then perform a critical external validation: we replicate the entire comparison on a second dataset, MMFi~\cite{yang2023mmfi} (Section~\ref{sec:boundary}). \textbf{The +8\,pp gain vanishes}---multi-scale is nominally ahead ($+1.84$\,pp vs.\ RMS-agg, $p = 9.2\times10^{-4}$; $+0.93$\,pp vs.\ multi-\Rx, $p = 0.116$, not significant; v1 path-labeled cache, $n = 1080$ samples). t-SNE visualization reveals the geometric reason: MMFi's four environments are indistinguishable in the \BVP{} feature space, meaning there is no discriminative signal for multi-\Rx{} to amplify. \textbf{Multi-\Rx{} is not unconditionally effective; it is a conditional improvement whose ceiling depends on the dataset's inherent separability.}

In summary, this paper makes four contributions:
\begin{itemize}
    \item \textbf{C1 (Methodology).} We establish a paired leave-one-domain-out $\times$ multi-seed $\times$ statistical-power-aware cross-domain benchmark protocol for \CSI{} gesture recognition, with both strict paired samples and broader unpaired multi-seed runs. The protocol, code, and all paired evaluations are released.
    \item \textbf{C2 (Empirical Finding).} We show that several recent algorithmic improvements in \CSI{} cross-domain recognition fail to reach significance ($|d| < 0.3$, $p > 0.28$), and that this failure is statistically meaningful---not a sample-size artifact---via \MDE{} analysis.
    \item \textbf{C3 (Mechanism).} We identify the \BVP{} multi-\Rx{} expansion as a stable improvement on Widar3.0 (+8.14\,pp, $d = 1.71$), and provide mechanistic evidence (antenna participation ratio, Fisher discriminability) that attributes the gain to front-end signal preservation rather than back-end algorithmic novelty.
    \item \textbf{C4 (Boundary).} We demonstrate, through cross-dataset replication and t-SNE visualization, that the multi-\Rx{} gain is conditional: it is bounded by the separability of the underlying feature space. When domains are indistinguishable (MMFi), multi-\Rx{} provides no measurable benefit and multi-scale is nominally optimal---the Widar advantage does not transfer (the residual nominal difference is not statistically significant on the v1 cache). \textbf{Updated C4 (v2 cache reanalysis, Section~\ref{sec:boundary}):} on the segment-level MMFi v2 cache ($n=16{,}447$), the BVP representation is in fact \emph{highly discriminative} for the 27-class task (between-vs-within variance ratio $= 617.85$), the LeNet-vs-CBAM+ResNet18 paired comparison becomes statistically indistinguishable ($\Delta = -0.33$\,pp, $p = 0.564$, 95\% CI spans zero, $n=3$ seeds), and the apparent representation flip in C7a--c is consistent with being a property of the v1 path-labeled cache construction rather than of MMFi itself. Dataset-conditioned optimality generalizes: on real-data benchmarks (Widar3.0 50k, MMFi v2 16{,}447) no backbone difference is statistically significant ($p > 0.05$; 95\,\% CIs span zero), the nominally significant 12k-subset result ($p = 0.028$) does not survive the 50k re-evaluation, and the MMFi v1 cache numbers are artifacts of the path-labeled cache construction.
\end{itemize}

% =============================================================================
\section{Related Work}
\label{sec:related}

We organize prior work into three streams most relevant to our contributions.

\subsection{Input-Side Reproducibility in \CSI{} Sensing}

Standardization protocols for \CSI{} inputs have been proposed~\cite{zhang2026sdp} but adoption remains uneven: a 2026 audit~\cite{guarino2026} reports that, in its sample, fewer than 35\% of recent \CSI{} papers release both their preprocessing pipeline and per-sample variance. A recent survey~\cite{wang2026survey} organises the \CSI{} processing pipeline into a four-stage taxonomy (experimental setup, signal preprocessing, feature learning, model deployment) but stops short of prescribing statistical evaluation. Most relevant to our paper are large-scale benchmarks that compare shallow and deep models~\cite{yang2023sensefi}, cross-domain adaptation methods~\cite{sterzinger2026datta}, and industrial-strength open-source deployments~\cite{cohen2026ruview}---all of which we treat as \emph{Reality Checks} for our protocol (Section~\ref{sec:intro}). None of these works, however, performs the \emph{paired multi-seed re-evaluation with statistical-power pre-registration} that we present.

Beyond benchmarking, the \CSI{} sensing community is rapidly adopting self-supervised and foundation-model paradigms. WiFi-JEPA~\cite{kim2026wifijepa} applies joint-embedding self-supervision to WiFi-CSI 3D human pose estimation, while AM-FM~\cite{zhu2026amfm} trains a foundation model for ambient-intelligence tasks through WiFi. We position such methods as \emph{representation-layer} counterparts to our \emph{evaluation-layer} protocol: they aim to learn transferable features, whereas our paired statistical instrument is what determines whether a reported gain is real rather than a seed/fold artifact.

\subsection{Statistical Methods for Deep Learning Benchmarks}

The broader machine learning community has long argued for variance reporting~\cite{henderson2018, bouthillier2021}, paired analysis~\cite{dror2018, colas2018}, and confidence intervals on effect sizes~\cite{bouthillier2021, dror2018}. We adopt three specific practices from this literature: (i)~\emph{paired-difference} tests rather than independent-sample tests~\cite{colas2018}; (ii)~\emph{Cohen's $d$ with Hedges' $g$ correction} for small samples~\cite{dror2018}; and (iii)~\emph{Minimum Detectable Effect (MDE)} pre-registration~\cite{dror2018, pineau2021}. To our knowledge, no \CSI{} cross-domain paper prior to ours has applied all three together.

\subsection{BVP and Multi-Antenna Systems}

The BVP operator~\cite{zheng2019widar} projects \CSI{} onto body-coordinate velocity features, producing a Doppler-time representation per antenna. WiTraj~\cite{wu2021} showed that aggregating across multiple receivers eliminates Doppler ambiguity. Wang \etal~\cite{wang2026survey} discuss multi-antenna systems but do not distinguish between \emph{aggregation} (collapse A $\to$ 1) and \emph{preservation} (retain A as a tensor axis). Our \MultiRx\ expansion falls in the latter category and provides what we believe to be the first principled comparison between these two strategies under a paired protocol.

\subsection{Neighboring Protocol and Reproducibility Efforts}
\label{sec:related-neighbors}

Three recent efforts intersect our contribution and require explicit positioning.
\emph{(i)~Standardized preprocessing (data layer).} Zhang~\etal~\cite{zhang2026sdp}
propose SDP, a unified protocol and benchmarking framework (with the open-source
\texttt{wsdp} backend) that standardizes \CSI{} input preprocessing and reportedly
reduces inter-seed variance by roughly $88\%$ on a 12,000-sample Widar3.0 subset (7.33\,\% of the full 163{,}650-sample unique link-level pool). \emph{Honest re-evaluation on a 50{,}000-sample subset (30.55\,\% of the full pool, paired t-test over 18 strict pairs) confirms the qualitative point that backbone ranking is dominated by protocol-induced uncertainty: the direction is stable---the smaller LeNetCSI\_Attn backbone wins on both subsets (12k: $\Delta = -4.61$\,pp for CBAM+ResNet18, $p = 0.028$; 50k: $\Delta = -5.75$\,pp, 95\,\% CI $[-11.77, +0.26]$, $p = 0.0609$, $d = -0.44$, 14/18 wins for LeNetCSI\_Attn)---but the significance does not survive the subset change (Section~\ref{sec:results}, Table~\ref{tab:main}, rows C10/C10v2). SDP standardizes the \emph{data layer} but cannot adjudicate which backbone wins; both call for paired statistical evaluation to obtain stable rankings.}
Our paired-LODO protocol operates at the \emph{evaluation layer} and is complementary: it can be run on
SDP-standardized inputs to yield comparable \emph{conclusions}, not merely
comparable \emph{data}.
\emph{(ii)~Protocol-induced uncertainty (measurement).} Da~Silva and Monahan~\cite{dasilva2026pss}
introduce a Protocol Sensitivity Score and a winner-swap test, showing that
between-protocol ranking displacement dwarfs within-protocol variance in PM10
forecasting. We treat our paper as the \emph{cross-domain instantiation} of this
thesis: we demonstrate the same protocol-induced uncertainty inside cross-domain
\CSI{} sensing and, crucially, add the cross-domain empirical evidence that LODO
exposes wins which single-domain cross-validation conceals.
\emph{(iii)~Uncertainty quantification for deployed models.} Li~\etal~\cite{li2026conformal}
apply conformal prediction to produce valid prediction sets for WiFi sensing models
in the wild. Their target is the \emph{per-sample predictive uncertainty of a trained
model}; ours is the \emph{meta-level trustworthiness of the comparison protocol
itself}. The two are orthogonal and composable: conformal prediction flags uncertain
test points while our protocol flags uncertain \emph{leaderboard claims}.
\emph{(iv)~Coordinate overfitting (mechanism).} Jia~\etal~\cite{chen2026perceptalign}
attribute cross-layout failure to \emph{coordinate overfitting} and report $>60\%$
cross-domain error reduction via geometry-aware alignment. Theirs is a
\emph{mechanism claim}; our paired-LODO is the \emph{evaluation instrument} that
lets such mechanism claims be compared fairly---without paired statistics, a reported
``$>60\%$ reduction'' may be partly a seed/fold artifact rather than a genuine effect.

% =============================================================================
\section{Methodology}
\label{sec:methodology}

Our methodology has five components: the \BVP{} operator (background), the multi-\Rx{} expansion (our modification), the \LODO\ evaluation protocol (cross-environment design), the paired multi-seed statistical protocol (variance control), and the \MDE\ framework (power pre-registration). Each is described below. Figure~\ref{fig:pipeline} gives an overview of the complete pipeline.

\begin{figure}[!t]
    \centering
    \includegraphics[width=\columnwidth]{figs/fig0_pipeline.png}
    \caption{Reproducibility audit pipeline. The flow proceeds through seven stages: raw \CSI\ input, \BVP\ preprocessing, representation choice (Multi-Rx, Multi-scale, or RMS-agg), backbone, \LODO\ protocol, paired $\times$ multi-seed evaluation, and statistical inference (paired $t$-test, Cohen's $d$, \MDE, t-SNE). The inset violin plot previews the dataset-dependent representation choice (Widar: Multi-Rx best; MMFi v1 cache: Multi-scale nominally best); see Section~\ref{sec:boundary}.}
    \label{fig:pipeline}
\end{figure}

\subsection{BVP Operator (Background)}

The Body-coordinate Velocity Profile (\BVP) operator~\cite{zheng2019widar} applies a Short-Time Fourier Transform along the packet dimension, removes the static component via Moving Target Indication (MTI), and projects the resulting Doppler spectrum onto body-coordinate axes via the linearly constrained minimum variance (LCMV) beamformer. For an input tensor $\mathbf{X} \in \mathbb{R}^{T \times A \times S}$ ($T$ packets, $A$ receive antennas, $S$ subcarriers), the BVP output is $\mathbf{B} \in \mathbb{R}^{T' \times S \times D}$ where $D$ is the Doppler-bin count. The canonical implementation aggregates across antennas before beamforming, producing a single tensor per packet.

\subsection{Multi-Rx Expansion (Our Modification)}

We depart from the canonical implementation by \emph{retaining} the antenna dimension throughout BVP computation. Concretely, for each antenna $a \in \{1, \ldots, A\}$ we compute a separate Doppler projection $\mathbf{B}_a \in \mathbb{R}^{T' \times S \times D}$, and stack them into a 5-D tensor $\mathbf{B}_{\text{multi}} \in \mathbb{R}^{N \times A \times S \times D \times T'}$. On Widar3.0, this yields a $(N, 9, 30, 33, 25)$ tensor; on MMFi, a $(N, 10, 114, 33, 29)$ tensor. The backbone (\textit{LeNetCSI\_Attn}) is applied independently to each antenna slice before feature fusion, so the network can learn per-antenna discriminative patterns. The full ablation chain is:
\begin{equation}
    \underbrace{\text{V1 (Multi-Rx)}}_{\text{keep } A=9}
    \xrightarrow{\text{antenna aggregation}}
    \underbrace{\text{V2 (RMS-agg)}}_{\text{collapse } A \to 1}
    \xrightarrow{\text{hop change}}
    \underbrace{\text{V3 (h16 RMS)}}_{\text{baseline}}
\end{equation}

\subsection{LODO Evaluation Protocol}

We adopt Leave-One-Domain-Out (\LODO) cross-validation: for $D$ domains, we train on $D-1$ domains and test on the held-out one, repeating $D$ times so every domain is the test set once. Widar3.0 has 6 user rooms (d1, d3, d5, d6, d7, d9); MMFi has 4 environments (E1--E4). The reported accuracy is the macro-average across all $D$ folds.

\subsection{Paired Multi-Seed Statistical Protocol}

Each (method, fold) pair is trained under $K = 5$ random seeds, yielding \emph{up to} $D \times K$ paired samples per comparison; the actual number of strict paired samples is determined by the overlap of seed sets between the method and its baseline. We report:
\begin{itemize}
    \item \textbf{Mean paired $\Delta$:} $\bar{\Delta} = \frac{1}{DK} \sum_{d,k} (\text{acc}_{\text{method}}(d,k) - \text{acc}_{\text{baseline}}(d,k))$, in percentage points (pp).
    \item \textbf{Paired $t$-statistic and $p$-value:} testing $H_0: \Delta = 0$ against $H_1: \Delta \neq 0$.
    \item \textbf{Cohen's $d$ with Hedges' $g$ correction:} $d = \bar{\Delta} / s_{\Delta}$ with $g = d \cdot (1 - 3/(4n - 1))$ for small-sample bias correction; 95\% CI via Hedges' approximation $\text{SE}_d = \sqrt{1/n + d^2 / (2n)}$.
    \item \textbf{95\% CI of $\bar{\Delta}$:} standard $t$-interval at $\alpha = 0.05$.
\end{itemize}

\subsection{Statistical Foundation: Closed-Form Bounds for the Paired Audit}
\label{sec:statfoundation}

We formalize the three statistical primitives that govern the paired audit. All bounds are derived from the standard Gaussian/Student-$t$ model and require only that the paired differences $\Delta_{d,k}$ have finite variance; the proofs are deferred to Appendix~\ref{app:proofs}.

\begin{theorem}[Paired Sample Size for Power $(1-\beta)$ at Effect $\varepsilon$]
\label{thm:paired_n}
Let $\Delta_1, \ldots, \Delta_n$ be paired differences with mean $\varepsilon$ and variance $\sigma_\text{within}^2$. The paired $t$-test rejects $H_0: \mathbb{E}[\Delta] = 0$ at significance $\alpha$ with asymptotic power $1-\beta$ when
\begin{equation}
    n \;\ge\; \frac{(z_{1-\alpha/2} + z_{1-\beta})^2 \, \sigma_\text{within}^2}{\varepsilon^2}\,,
    \label{eq:paired_n}
\end{equation}
where $z_p$ is the standard normal quantile.
\end{theorem}

\noindent\textbf{Substituting our empirical regime.} Widar3.0 paired runs yield $\sigma_\text{within} \approx 4.8$\,pp (Section~\ref{sec:variance}). To detect $\varepsilon = 2.0$\,pp at $\alpha = \beta = 0.05$ requires $n \ge (1.96 + 0.84)^2 \cdot 23.04 / 4.00 \approx 56$ paired samples --- more than triple our actual $n \le 18$. Consequently, single-seed reports of $\le 2$\,pp gain are \textbf{provably underpowered} under any paired design at our sample budget, which explains the empirical failure rate in Section~\ref{sec:results}.

\begin{theorem}[MDE Closed Form]
\label{thm:mde_closed}
At significance $\alpha$ and power $1-\beta$, the Minimum Detectable Effect (\MDE{}) for a paired $t$-test on $n$ samples of within-pair variance $\sigma_\text{within}^2$ is
\begin{equation}
    \MDE{} = (t_{\alpha/2,\,n-1} + t_{\beta,\,n-1}) \cdot \frac{\sigma_\text{within}}{\sqrt{n}}\,.
    \label{eq:mde_closed_eq}
\end{equation}
For $\alpha = \beta = 0.05$ and $n \ge 10$, this approximates $2.91 \cdot \sigma_\text{within}/\sqrt{n}$ with relative error below 4\%.
\end{theorem}

\begin{theorem}[Hedges' $g$ Small-Sample Correction]
\label{thm:hedges}
For Cohen's $d = \bar{\Delta}/s_\Delta$ computed from $n$ paired differences, the unbiased estimator of the population effect size is
\begin{equation}
    g \;=\; d \cdot \left( 1 - \frac{3}{4n - 1} \right)\,,
    \qquad \mathrm{SE}_g \;=\; \sqrt{\frac{1}{n} + \frac{g^2}{2n}}\,.
    \label{eq:hedges_eq}
\end{equation}
\end{theorem}

\begin{corollary}[Reproducibility Audit Criterion]
\label{cor:audit_criterion}
A reported effect is declared \emph{real} only when all three conditions hold simultaneously: (i) $|\bar{\Delta}| > \MDE{}$, (ii) $p < \alpha$, and (iii) $|g| \ge 0.2$. Single-criterion acceptance inflates the false-positive rate under realistic $\sigma_\text{within}$ budgets: criterion (i) alone is necessary but not sufficient when the effect direction is unconstrained, and criterion (iii) alone ignores the sampling-distribution width of the underlying estimator.
\end{corollary}

\noindent\textbf{Why this matters for CSI sensing.} The closed-form bounds above reduce the audit to a pre-registration exercise: given a target effect $\varepsilon$ and the empirically-observed $\sigma_\text{within}$, the reviewer can compute the required $n$ before training a single model, and flag any single-seed report whose claimed gain lies below the resulting \MDE{} as \emph{provably underpowered}. We deploy this audit uniformly across the 12 comparisons in Section~\ref{sec:results} and record the per-comparison MDE in Appendix~\ref{app:mde}.

\subsection{MDE Framework (Power Pre-Registration)}
\label{sec:mde_framework}

Before running any comparison we pre-register the Minimum Detectable Effect (\MDE) at power $1 - \beta = 0.80$ and significance $\alpha = 0.05$:
\begin{equation}
    \MDE{} = (t_{\alpha/2,\,n-1} + t_{\beta,\,n-1}) \cdot \frac{\sigma_{\text{within}}}{\sqrt{n}}
        \approx 2.91 \cdot \frac{\sigma_{\text{within}}}{\sqrt{n}}
    \label{eq:mde}
\end{equation}
For our paired design (Widar: up to $D \cdot K = 6 \cdot 5 = 30$ runs per comparison, of which 18 have a strict paired baseline in the 3-shared-seed runs and 30 for the C2/C5 5-seed sets; MMFi: $n = 4 \cdot 5 = 20$), MDE reduces to $0.53 \cdot \sigma_{\text{within}}$. Empirically, $\sigma_{\text{within}}$ varies from 1.0\,pp (RMS-agg) to 8.2\,pp (Multi-Rx), yielding MDEs of 0.5\,pp to 4.4\,pp. We declare a comparison ``real'' only when $|\bar{\Delta}| > \MDE$ and $p < 0.05$ and $|g| \geq 0.2$ \emph{simultaneously}; otherwise it is consistent with noise given our measurement precision.

% =============================================================================
\section{Experimental Setup}
\label{sec:setup}

\subsection{Datasets}

\textbf{Widar3.0}~\cite{zheng2019widar}: the official release comprises 17{,}000 gesture instances---12{,}000 hand-gesture instances covering six gestures (push\&pull, sweep, clap, slide, draw-circle, draw-zigzag; $16 \text{users} \times 5 \text{positions} \times 5 \text{orientations} \times 6 \text{gestures} \times 5 \text{instances}$) across three environments, plus 5{,}000 draw-number instances from two volunteers. Each gesture instance is captured by multiple Tx--Rx links; after deduplication our pool contains 163{,}650 unique link-level records. In the released archives, gesture indices 1--5 appear consistently across 13 domain archives, indices 6--8 occur only in two additional archive sessions, and index~9 is absent; the released index set, rather than any nominal taxonomy, is what our caches inherit. We hold out six domain archives (d1, d3, d5, d6, d7, d9) in turn as 6 folds for \LODO. Each sample is a $(T{=}256, A{=}9, S{=}30)$ complex-valued \CSI{} tensor at 1000\,Hz packet rate.

\textbf{MMFi}~\cite{yang2023mmfi}: 27 gestures (Table~\ref{tab:main} C7a-c uses the v1 path-labeled cache across all 27 classes; Section~\ref{sec:boundary} re-runs the C10 capacity control on the v2 segment-level cache with 16,447 segments across all 27 gestures, while the C7a--c representation rows remain on the v1 cache and are flagged as v1-cache-conditioned), 4 environments (E1--E4 with different rooms and orientations), 40 users; \CSI{} is sampled at 100\,Hz on a 10-antenna array across 114 subcarriers.

\subsection{Backbone}

All methods share the \textit{LeNetCSI\_Attn} backbone: three convolutional blocks (32/64/128 channels) followed by a 4-head attention pooling layer and an MLP head. The backbone is intentionally simple so that gains cannot be attributed to architectural novelty.

\subsection{Methods Compared}

We compare ten method axes spanning three categories (yielding 12 paired comparisons; Table~\ref{tab:main} reports 13 rows counting the C10v2 re-evaluation of the capacity control on a second subset). \textbf{Architectural (attention variants, 3 variants):} LeNet (baseline), LeNet+Attn vs.\ LeNet+Attn+MHA (C1). \textbf{Algorithmic (4 axes):} DANN (GRL, C4), Snapshot$\times$3 ensemble (C3), Mixup (C8), Reweighting (C9). \textbf{Representational (3 axes, 5 variants):} Widar-orig (RMS hop=16, C6 baseline), Multi-scale (hop $\in \{8, 16, 32\}$, C2 / C7b / C7c), Multi-Rx (keep $A=9$, ours, C5 / C6 / C7a / C7c), RMS-agg (vanilla, C5 / C7a / C7b), Single-Rx ($A=1$, reference). Each axis runs up to $K = 5$ seeds $\times$ $D$ folds = 30 runs on Widar, 20 on MMFi, with the number of strict paired samples determined by the baseline's seed coverage; the exact $n$ per comparison is reported in Appendix~C. As a \emph{larger-backbone capacity control}, we additionally evaluate a CBAM+ResNet18 backbone (11.42M parameters, roughly $89\times$ the size of LeNetCSI\_Attn) under the same protocol; this comparison is reported as \textbf{C10} and serves as a sanity check that the observed effects are not driven by model capacity alone. The original C10 evaluation used a 12{,}000-sample subset of the dedup-ed Widar3.0 pool (7.33\,\% of the full 163{,}650 unique link-level samples); to check the generality of the conclusion we additionally evaluated \textbf{C10v2} on a 50{,}000-sample subset (30.55\,\% of the full pool) under the identical paired LODO protocol. Both backbones consume the identical folds, seeds, and labels within each subset, so the capacity contrast is internal to this fixed labeling (the released index set is disclosed in Section~\ref{sec:setup}).

\subsection{Statistical Pipeline}

All numbers reported in Section~\ref{sec:results} are computed by an automated pipeline (released with the code) that: (i)~loads the raw per-seed accuracies; (ii)~computes paired differences across each (fold, seed) pair; (iii)~runs a paired $t$-test against zero; (iv)~computes Cohen's $d$ with Hedges' $g$ correction and its 95\% CI; (v)~runs an MDE check at $\alpha = 0.05$, power $= 0.80$; (vi)~flags ``non-significant'' vs ``real'' per Eq.~\ref{eq:mde}.

% =============================================================================
\section{Main Results: Twelve Paired Comparisons Under the Paired Audit}
\label{sec:results}

\begin{figure*}[!t]
    \centering
    \includegraphics[width=0.92\textwidth]{figs/fig1_main_boxplot.png}
    \caption{Paired forest plot of the 9 pre-registered paired comparisons (dual-panel; the later additions C8--C10v2 are reported in Table~\ref{tab:main}). \textbf{Left:} mean $\Delta$ with 95\% CI of the absolute gain (pp). \textbf{Right:} Cohen's $d$ with Hedges' $g$-correction and 95\% CI, with $|d| \in \{0.2, 0.5, 0.8\}$ reference thresholds for small/medium/large effects. Four comparisons reach $p < 0.05$ with positive $\Delta$ (C5, C6, C7a, C7b); three satisfy all three audit criteria of Corollary~\ref{cor:audit_criterion} (C5, C6, C7b). Across the full audit, five comparisons cross $|g| \geq 0.8$: C5 ($g = 0.94$), C6 ($g = 1.68$), C7b ($g = 0.84$), and the two harmful controls C8 ($g = -0.90$), C9 ($g = -1.11$). All four ``algorithmic'' methods (C1, C3, C4, C2) fail both tests and are consistent with noise. C8, C9 demonstrate that strong algorithmic interventions can \emph{hurt} multi-Rx performance under paired evaluation; the C10 backbone control points in the same (negative) direction, though its significance is subset-fragile (Section~\ref{sec:results}).}
    \label{fig:forest}
\end{figure*}

We systematically evaluated ten recent improvements under a paired LODO protocol on Widar3.0 (6 folds; the C10 capacity control uses $n = 18$ strict paired samples from 3 shared seeds; other Widar comparisons use $n = 18$ or $n = 30$ strict paired samples depending on the seed-coverage of the baseline---see Appendix~C for the exact $n$ per comparison). Figure~\ref{fig:forest} presents the results as a dual-panel paired forest plot; Table~\ref{tab:main} reports the numerical values. Of the ten methods, three representation-family comparisons reach significance with $|g| \geq 0.8$ (C6 $g = 1.68$, C5 $g = 0.94$, C7b $g = 0.84$); C7a (multi-\Rx\ vs RMS-agg on MMFi) reaches $p = 0.020$ with a medium effect ($g = 0.55$) whose $|\bar{\Delta}|$ falls below its pre-registered \MDE. \textbf{C10 (the larger-backbone capacity control, CBAM+ResNet18) reaches $\Delta = -4.61$\,pp ($p = 0.028$, $d = -0.52$, $3/18$ wins for CBAM+ResNet18) \emph{on the original 12k Widar3.0 subset} (7.33\,\% of the full unique pool), suggesting that the larger backbone may slightly hurt cross-domain accuracy. Honest re-evaluation \textbf{C10v2} on a 50{,}000-sample subset (30.55\,\% of the full pool; same protocol, 18 strict pairs) returns $\Delta = -5.75$\,pp (95\,\% CI $[-11.77, +0.26]$, $p = 0.0609$, $d = -0.44$, 4/18 wins for CBAM+ResNet18).} The direction is therefore \emph{stable} across subsets---the smaller backbone wins on both---but the \emph{significance does not survive}: the nominally significant 12k result ($p = 0.028$) degrades to $p = 0.0609$ with a CI spanning zero at 50k, and the paired-difference standard deviation inflates from 8.9 to 13.0\,pp. Both subsets are drawn from the same dedup-ed release pool, whose gesture indices 1--5 are consistent across 13 domain archives while indices 6--8 (${\sim}8.8\,\%$ of the 50k subset) come from only two archive sessions; because both backbones consume identical folds, seeds, and labels, the paired capacity contrast is internal to this fixed labeling. We report both numbers transparently: this is evidence that the \emph{statistical significance} of a backbone comparison is subset-fragile under protocol-induced variance, not evidence of a winner swap---and we accordingly retract any reading of ``which backbone wins is a property of the evaluation subset.'' We also remind the reader that the 50k $p = 0.0609$ marginally exceeds the $\alpha = 0.05$ threshold, so the 50k result is consistent with both a small negative and a near-null effect; we do not claim the larger backbone is universally preferable or harmful. The remaining six comparisons all fail both the $p < 0.05$ test and the $|g| \geq 0.2$ test; we report them as \emph{non-significant} under the present protocol rather than as ``artifacts''.

\begin{table}[!t]
    \centering
    \caption{Main result: 12 primary paired comparisons across Widar and MMFi, plus the C10v2 re-evaluation (13 rows). $\Delta$ in pp; $d$ = Cohen's $d$; $g$ = Hedges' $g$-corrected; CI at 95\%; $n$ = number of strict paired samples. Bold rows cross $|g| \geq 0.8$ or are pre-specified ``capacity/hurts'' controls. The C10 capacity control was originally on a 12k subset (7.33\% of full pool) and re-evaluated as \textbf{C10v2} on a 50k subset (30.55\% of full pool); both rows are reported. Both Widar subsets are drawn from the same dedup-ed release pool (163{,}650 unique link-level records; released gesture indices 1--8, with indices 6--8 confined to two archive sessions), and the backbone contrast is internal to this fixed labeling. Other Widar comparisons use $n = 18$ or $n = 30$ strict paired samples (Appendix~C); MMFi uses $n = 20$ (4 folds $\times$ 5 seeds). \textbf{Cache note:} C7a-c are computed on the v1 path-labeled MMFi cache (1080 samples, 27$\times$40 artificial balance) as documented in Section~\ref{sec:boundary}; Section~\ref{sec:boundary} additionally re-runs the C10 capacity control on the v2 segment-level MMFi cache (16,447 segments from \texttt{MMFi\_action\_segments.csv}), and reports the LeNet-vs-CBAM+ResNet18 paired comparison on the v2 cache as $\Delta = -0.33$\,pp, $p = 0.564$, 95\% CI spans zero ($n = 3$ seeds)---i.e., the v2 backbone comparison is statistically indistinguishable. The v1 cache numbers in C7a-c are retained for transparency but should be read as v1-cache-conditioned rather than as a representation-level conclusion about MMFi.}
    \label{tab:main}
    \scriptsize
    \setlength{\tabcolsep}{2pt}
    \begin{tabular}{@{}l >{\raggedright\arraybackslash}p{3.6cm} l c c c c >{\raggedright\arraybackslash}p{1.5cm}@{}}
    \toprule
    ID & Comparison & Dataset & $\Delta$ (pp) & $d$ & $g$ & $p$ & Sig. \\
    \midrule
    C1 & MHA vs Conv.\ attn.               & Widar & $+1.13$   & $+0.17$ & $+0.16$ & 0.478 & ---  \\
    C2 & Multi-scale vs Multi-\Rx          & Widar & $-1.28$   & $-0.14$ & $-0.14$ & 0.450 & ---  \\
    C3 & Snapshot$\times$3 vs single       & Widar & $+0.23$   & $+0.15$ & $+0.14$ & 0.537 & ---  \\
    C4 & DANN vs source-only               & Widar & $-1.22$   & $-0.15$ & $-0.14$ & 0.530 & ---  \\
    \textbf{C5} & \textbf{Multi-\Rx\ vs RMS-agg} & \textbf{Widar} & $\mathbf{+7.69}$ & $\mathbf{+0.97}$ & $\mathbf{+0.94}$ & $\mathbf{1.1\!\times\!10^{-5}}$ & \textbf{yes} \\
    \textbf{C6} & \textbf{Multi-\Rx\ vs Widar-orig} & \textbf{Widar} & $\mathbf{+8.14}$ & $\mathbf{+1.71}$ & $\mathbf{+1.68}$ & $\mathbf{1.3\!\times\!10^{-6}}$ & \textbf{yes} \\
    C7a & Multi-\Rx\ vs RMS-agg (MMFi)     & MMFi & $+0.91$   & $+0.57$ & $+0.55$ & 0.020 & below \MDE \\
    \textbf{C7b} & \textbf{Multi-scale vs RMS-agg (MMFi)} & \textbf{MMFi} & $\mathbf{+1.84}$ & $\mathbf{+0.88}$ & $\mathbf{+0.84}$ & $\mathbf{9.2\!\times\!10^{-4}}$ & \textbf{yes} \\
    C7c & Multi-scale vs Multi-\Rx\ (MMFi)     & MMFi & $+0.93$   & $+0.37$ & $+0.35$ & 0.116 & --- \\
    \textbf{C8} & \textbf{Mixup vs vanilla Multi-\Rx} & \textbf{Widar} & $\mathbf{-7.90}$ & $\mathbf{-0.96}$ & $\mathbf{-0.90}$ & $\mathbf{0.0066}$ & \textbf{hurts} \\
    \textbf{C9} & \textbf{Reweight vs vanilla Multi-\Rx} & \textbf{Widar} & $\mathbf{-10.17}$ & $\mathbf{-1.19}$ & $\mathbf{-1.11}$ & $\mathbf{0.0017}$ & \textbf{hurts} \\
    \textbf{C10} & \textbf{CBAM+ResNet18 vs Attn (cap, 12k)} & \textbf{Widar} & $\mathbf{-4.61}$ & $\mathbf{-0.52}$ & $\mathbf{-0.49}$ & $\mathbf{0.028}$ & \textbf{non-positive} \\
    \textbf{C10v2} & \textbf{CBAM+ResNet18 vs Attn (cap, 50k)} & \textbf{Widar} & $\mathbf{-5.75}$ & $\mathbf{-0.44}$ & $\mathbf{-0.42}$ & $\mathbf{0.061}$ & \textbf{consistent (n.s.)} \\
    \bottomrule
    \end{tabular}
\end{table}

\paragraph{Algorithmic methods (C1, C3, C4, C2)---non-significant under paired evaluation.}
The four algorithmic methods (MHA, Snapshot$\times$3, DANN, multi-scale vs multi-\Rx\ on Widar) all produce $|\bar{\Delta}| \leq 1.28$\,pp with $p > 0.45$. Under our 18-sample paired designs these effects fall below their pre-registered \MDE{} values (1.10--5.65\,pp; Appendix~C). The published gains of 2--5\,pp claimed in the original papers are therefore not distinguishable from noise in our paired evaluation; we do not claim the original results are incorrect.

\paragraph{Representational methods (C5, C6, C7a--c)---the representation axis.}
The \MultiRx\ expansion produces $\Delta = +7.69$\,pp vs RMS-aggregation ($g = 0.94$, $p = 1.1 \times 10^{-5}$) and $\Delta = +8.14$\,pp vs the canonical Widar-orig ($g = 1.68$, $p = 1.3 \times 10^{-6}$). On MMFi (v1 path-labeled cache) the multi-\Rx\ advantage over RMS-agg shrinks to $+0.91$\,pp ($p = 0.020$, below its \MDE), multi-scale beats RMS-agg by $+1.84$\,pp ($p = 9.2 \times 10^{-4}$), and multi-scale's nominal edge over multi-\Rx\ is not significant ($+0.93$\,pp, $p = 0.116$): the Widar representation advantage does not transfer to MMFi. The dataset-optimal representation is therefore \emph{dataset-dependent} on the v1 cache---a characterization revisited under the v2 reanalysis in Section~\ref{sec:boundary}.

% =============================================================================
\section{Variance Decomposition: Why Seeds Mislead}
\label{sec:variance}

We decompose the total variance of the per-fold accuracy into four components: \emph{seed} (between-seed), \emph{fold} (between-domain, treated as fixed), \emph{variant} (between-representation), and \emph{fold-mean residual} (everything else, including sample stochasticity within a fold-seed pair).

\begin{figure}[!t]
    \centering
    \includegraphics[width=\columnwidth]{figs/fig2_variance_decomp.png}
    \caption{Variance decomposition for the 4 BVP variants (multi-Rx, multi-scale, RMS-agg, single-Rx) on Widar3.0 LODO, $K=5$ seeds, $D=6$ folds. The dominant component is \emph{seed} (4.8\,pp on average across the four variants), which is 4.7$\times$ larger for Multi-Rx than for the other variants.}
    \label{fig:variance}
\end{figure}

\textbf{Finding.} The seed component dominates: $\sigma_{\text{seed}} = 4.8$\,pp on average (Fig.~\ref{fig:variance}), versus $\sigma_{\text{fold}} = 1.4$\,pp and $\sigma_{\text{residual}} = 0.7$\,pp. Crucially, the seed component is \emph{5$\times$ larger for Multi-Rx than for RMS-agg}, even though both are evaluated under the same protocol. This means that the more expressive representation is also the \emph{more variable} one---a finding that single-seed studies are structurally unable to detect.

% =============================================================================
\section{MDE Analysis: Power in Practice}
\label{sec:mde}

We pre-register the Minimum Detectable Effect (\MDE) at $\alpha = 0.05$, power $= 0.80$ for each comparison before running it (Eq.~\ref{eq:mde}). Figure~\ref{fig:mde} shows the MDE self-validation.

\begin{figure}[!t]
    \centering
    \includegraphics[width=\columnwidth]{figs/fig3_mde_validation.png}
    \caption{\MDE\ self-validation. For each of the 9 pre-registered paired comparisons (C1--C7c), the bar shows the observed $|\bar{\Delta}|$; the annotation reports the \MDE\ given the observed within-comparison $\sigma_{\text{within}}$ and $n$ (Eq.~\ref{eq:mde}). Six of the nine comparisons fall below their pre-registered \MDE{} (not distinguishable from noise at our precision); C5, C6, and C7b cross it.}
    \label{fig:mde}
\end{figure}

\textbf{Finding.} The four algorithmic methods (C1, C2, C3, C4) all produce $|\bar{\Delta}|$ below their respective pre-registered \MDE{} values ($|\Delta| \le 1.28$\,pp vs \MDE{} $\in [1.10, 5.65]$\,pp), so under the audit criterion of Corollary~\ref{cor:audit_criterion} they are declared \emph{not distinguishable from noise} at our measurement precision. Only the representational methods C5 and C6 cross the \MDE\ bar by a wide margin. We emphasize that ``below \MDE'' is a statement about our detection threshold, not a proof of a null effect.

% =============================================================================
\section{Mechanism: Why Multi-Rx Helps}
\label{sec:mechanism}

We investigate the mechanism behind the +8.14\,pp gain via three quantitative lenses (Figure~\ref{fig:mechanism}).

\begin{figure*}[!t]
    \centering
    \includegraphics[width=0.92\textwidth]{figs/fig5_mechanism.png}
    \caption{Multi-\Rx\ mechanism decomposition (4-panel). \textbf{(a)} Antenna Participation Ratio (PR) distribution: PR $= 1.89 \pm 0.17$ (5\%--95\%: $[1.75, 2.24]$, $n=2500$ positions), meaning that on average only $\sim$2 of the 9 antennas are linearly independent. \textbf{(b)} Fisher $F$ per axis: Multi-\Rx's per-antenna $F$ is \emph{lower} than RMS-agg's per-axis $F$ (4.65 vs 23.05), but (c) information density is 5$\times$ higher. \textbf{(d)} Verdict: the gain comes from preserving density, not from increasing per-position discriminative power.}
    \label{fig:mechanism}
\end{figure*}

\subsection{Antenna Participation Ratio (PR)}

We define the antenna participation ratio of a per-position covariance matrix $\mathbf{C} \in \mathbb{R}^{A \times A}$ as
\begin{equation}
    \PR(\mathbf{C}) \;=\; \frac{\left(\sum_{i=1}^{A} \lambda_i\right)^2}{\sum_{i=1}^{A} \lambda_i^2},
    \label{eq:pr}
\end{equation}
where $\lambda_1 \geq \cdots \geq \lambda_A \geq 0$ are the eigenvalues of $\mathbf{C}$ (Roy \& Vetterli 2007). Empirically, $\PR = 1.89 \pm 0.17$ on Widar3.0, meaning that of the 9 antennas only $\sim 2$ are linearly independent. RMS-aggregation contracts all 9 antennas to a single slot, so it can encode at most $\PR = 1$ effective antenna; \MultiRx\ retains all 9 slots and thus can encode up to $\PR = 1.89$, an information-content gain of $1.89 / 1.00 = 1.89\times$ in the antenna dimension alone.

\subsection{Fisher $F$ Per Axis}

For each axis position $k$, we compute the Fisher criterion $F(k) = \sigma^2_{\text{between}}(k) / \sigma^2_{\text{within}}(k)$. On Multi-Rx the per-antenna $F_{\text{mean}} = 4.65$, while on RMS-agg $F_{\text{mean}} = 23.05$ (4.95$\times$ higher per slot). However, Multi-Rx retains $9\times$ more slots, so its \emph{total} Fisher information $\sum_k F(k)$ exceeds RMS-agg's by $1.82 \times$. RMS-aggregation \emph{loses} low-$F$ but high-density information that the backbone can still exploit.

\subsection{Information Density Fold}

The total effect decomposes as: $\rho_{\text{antenna}} = 0.20$ (5$\times$ density gain from antenna preservation) $\times \rho_{\text{hop}} = 1.00$ (no density change from hop) $= 0.20$ (V1 vs V3 total). The 5$\times$ density gain translates directly into the observed $\sim$8\,pp accuracy gain.

\subsection{Caveat: Per-Domain PR Does Not Predict Per-Domain Gain}

Plotting per-domain PR against per-domain Multi-Rx gain yields Spearman $\rho = +0.10$ (see Appendix B, \ref{fig:pr_gain}). The mechanism is \emph{global}, not per-domain: the participation ratio captures the average antenna redundancy, but the per-domain gain depends on the specific domain geometry that we explore next.

% =============================================================================
\section{Boundary: Why Multi-Rx Fails on MMFi}
\label{sec:boundary}

To test the generalizability of the Widar3.0 finding, we replicate the entire protocol on MMFi (27 classes, 4 environments).

\subsection{Dataset-Flip Discovery}

On MMFi the dataset-optimal representation is \emph{nominally} multi-scale rather than multi-\Rx: C7b multi-scale vs RMS-agg $\Delta = +1.84$\,pp ($p = 9.2\times10^{-4}$, above its pre-registered \MDE) and C7c multi-scale vs multi-\Rx\ $\Delta = +0.93$\,pp ($p = 0.116$, not significant; v1 path-labeled cache, $n = 20$ paired runs). The +8.14\,pp gain from multi-\Rx\ on Widar3.0 does not transfer.

\subsection{Geometric Explanation via t-SNE}

\begin{figure*}[!t]
    \centering
    \begin{subfigure}{0.48\textwidth}
        \includegraphics[width=\textwidth]{figs/fig_tsne_widar.png}
        \caption{Widar3.0 t-SNE: domains form partially separable clusters.}
    \end{subfigure}
    \begin{subfigure}{0.48\textwidth}
        \includegraphics[width=\textwidth]{figs/fig_tsne_mmfi.png}
        \caption{MMFi t-SNE: domains are \emph{indistinguishable}.}
    \end{subfigure}
    \caption{Domain separability in the \BVP\ feature space. On Widar3.0, the 6 domains form partially separable clusters, leaving room for multi-\Rx\ to exploit. On MMFi, the 4 domains are visually indistinguishable---there is no separability for multi-\Rx\ to amplify, and multi-scale (which captures gestural intra-class structure) becomes optimal.}
    \label{fig:tsne}
\end{figure*}

\textbf{Finding.} On Widar3.0 (Fig.~\ref{fig:tsne}\,a), the 6 domains form partially separable clusters in the \BVP\ feature space, leaving headroom for a representation that retains per-antenna discriminative power. On MMFi (Fig.~\ref{fig:tsne}\,b), the 4 domains are visually indistinguishable, so multi-\Rx\ has nothing to amplify; multi-scale, which captures gestural intra-class structure within each domain, becomes the dataset-optimal choice instead. \textbf{No single \BVP\ representation is universally optimal; the choice must be matched to the dataset's geometric structure.}

\subsection{Honest Limitation: Time Dimension}

A reviewer might ask whether the multi-Rx gain survives under temporal distribution shift. We do not have time-indexed \CSI\ data and so cannot answer this question directly; we flag it as a limitation and suggest it as future work. Independent evidence from a recent longitudinal study~\cite{vassallo2025time} suggests that \CSI\ features drift measurably over months, but no audit has yet been performed under paired multi-seed temporal evaluation.

\subsection{Honest Caveat: The v1 MMFi Cache Was Path-Labeled}

\textbf{Why we are re-running the MMFi benchmark.} After the initial submission we identified that the original 1080-case MMFi cache (used for the C7a--c numbers above) was constructed by taking the path-derived label (the \texttt{Axx} directory name) as the ground truth and treating the entire 297-frame \CSI\ sequence as one sample; this produced an \emph{artificial 27\,$\times$\,40 balanced} matrix (1080 samples) but did \emph{not} correspond to the action-segment-level ground truth shipped with the dataset. The MMFi repository provides a per-segment annotation file (\texttt{MMFi\_action\_segments.csv}) that defines the actual start/end frame for each gesture within a recording; the cache was constructed \emph{before} we had read this file, so the original numbers in Table~\ref{tab:main} (C7a\,$+0.91$\,pp, C7b\,$+1.84$\,pp, C7c\,$+0.93$\,pp) are valid only \emph{for that specific path-labeled subsample}. We retain those rows in Table~\ref{tab:main} for transparency but flag them as v1-cache-conditioned, not as a representation-level conclusion about the dataset.

\subsection{Real-Segment Reanalysis: MMFi v2 Cache}

\textbf{v2 cache construction.} Using \texttt{MMFi\_action\_segments.csv} as the ground truth and aggregating per-segment magnitude features across all 10 receiving antennas at 100\,Hz with a 96-frame stride, we obtain \textbf{16{,}447 segments} distributed across 27 classes with severe natural imbalance (A01 = 1547 vs A02 = 356; max/min = 4.35$\times$). The BVP representation on the v2 cache has shape $(16447, 10, 114, 33, 29)$ and is cached on disk as \texttt{\_mmfi\_v2\_bvp.h5} (72\,GB) with chunked streaming to fit the 67\,GB working set within our 68\,GB RAM budget.

\textbf{F-50 v2 signal-presence audit (label integrity, variance ratio, granularity).} On the v2 cache we re-ran the F-50 audit (label-integrity, between-vs-within variance, 5-class coarse retraining, raw-vs-BVP input compare). The label-integrity check found \textbf{0 issues} across all 27 classes $\times$ 4 environments. The between-vs-within variance ratio on the BVP representation is \textbf{617.85} (>>\,1), which by the standard interpretation rule rules out the ``representation at fault'' diagnosis the v1 numbers appeared to support. Coarse 5-class retraining reached \textbf{test = 0.318} (best-val 0.382; early stop ep11), compared to a 1/5 random baseline of 0.200. Granularity sweep on BVP (LeNet, E2 LODO) reached test = 0.603 on the 2-class grouping (daily vs rehab), 0.440 on 4-class, 0.319 on 5-class, 0.415 on 8-class, 0.307 on 12-class, and 0.169 on 27-class---confirming that on the real segment distribution, the bottleneck is \emph{class granularity}, not the BVP representation itself.

\textbf{F-48 v2 paired benchmark (3 seeds, LeNet vs CBAMResNet18).} We re-ran the C10 capacity control on the v2 cache under the same paired LODO protocol used for the Widar3.0 C10v2 evaluation: 3 seeds $\times$ 2 backbones (LeNetCSI\_Attn, 0.13M params; CBAM+ResNet18, 11.43M params) $\times$ E2 LODO (train=12513, val=786, test=3148) $\times$ multirx10 BVP input $\times$ best-val-state convention. Results: LeNet = 0.225, 0.228, 0.232 (mean 0.2284); CBAM = 0.229, 0.239, 0.227 (mean 0.2317); per-seed $\Delta$(CBAM$-$LeNet) = $+0.0041$, $+0.0111$, $-0.0054$ (mean $+0.33$\,pp). \textbf{Paired $t$-test: $t = 0.686$, $p = 0.5637$, Cohen's $d = +0.40$, 95\% CI $[-0.017, +0.024]$ (spans zero).} With only $n = 3$ seeds on a single E2 fold, however, the comparison is underpowered (its implied \MDE{} ${\approx}\,1.4$\,pp exceeds $|\Delta|$), so we report it as \emph{directionally near-null but statistically unresolved}, not as a proven null. On the v1 path-labeled cache the same two backbones are both near the 1/27 random baseline (LeNetCSI\_Attn 6.22\,\%, CBAM+ResNet18 4.56\,\%, $\Delta = -1.67$\,pp, $n = 20$): the v1 numbers are a cache artifact, not a model property.

\subsection{Updated Dataset-Dependent Optimality (Four Data Points)}

\begin{table}[!t]
    \centering
    \caption{Dataset-dependent backbone optimality: $\Delta$(CBAM+ResNet18 $-$ LeNetCSI\_Attn) on four (dataset, subset) combinations. The direction is stable across the two Widar3.0 subsets (the smaller backbone wins on both), but the nominally significant 12k result ($p = 0.028$) does not survive the 50k re-evaluation ($p = 0.061$, CI spans zero); on the MMFi v1 path-labeled cache both backbones sit near the 1/27 random baseline, and the v2 comparison is underpowered ($n{=}3$). The pattern supports \emph{protocol-induced uncertainty}: backbone-comparison significance is subset-fragile, and no backbone difference is statistically significant on any real-data benchmark.}
    \label{tab:ddopt}
    \scriptsize
    \setlength{\tabcolsep}{4pt}
    \begin{tabular}{lccccc}
    \toprule
    Subset & $n$ & Cache type & $\Delta$(pp) & $p$ & Sig. \\
    \midrule
    MMFi v1 (path-labeled)       & 1080  & fake balance & $-1.67$    & 0.001  & \emph{both $\approx$ random} \\
    Widar3.0 12k (7.33\% pool)   & 12000 & subsample    & $-4.61$    & 0.028  & \emph{nominally significant} \\
    Widar3.0 50k (30.55\% pool)  & 50000 & real         & $-5.75$    & 0.061  & \emph{n.s., CI spans 0} \\
    \textbf{MMFi v2 (real segs.)} & \textbf{16447} & \textbf{real} & $\mathbf{+0.33}$ & \textbf{0.564} & \textbf{underpowered ($n{=}3$)} \\
    \bottomrule
    \end{tabular}
\end{table}

\textbf{Updated claim.} The central claim of Section~\ref{sec:boundary} survives the v2 cache reanalysis but its mechanism becomes \emph{dataset-conditioned} rather than \emph{representation-flipped}. Across the real-data evaluations, no backbone difference is statistically significant (Widar3.0 12k: $\Delta = -4.61$\,pp, $p = 0.028$, nominally significant but not reproduced at 50k with $\Delta = -5.75$\,pp, $p = 0.061$, CI spanning zero; MMFi v2: $\Delta = +0.33$\,pp, $p = 0.564$, $n{=}3$ underpowered). The residual dataset-flip signal in C7a--c is a property of the v1 cache construction, not of MMFi itself; on real MMFi segments, BVP is highly discriminative (variance ratio = 617.85) and the bottleneck is class granularity (27 fine-grained gestures), not the representation choice. \textbf{We therefore revise C4 as follows: no single backbone or BVP representation is unconditionally optimal; apparent backbone or representation advantages that vanish under cache correction (MMFi v1 $\to$ v2) or whose significance does not survive subsampling (Widar3.0 12k $\to$ 50k) must be reported with their cache/subset construction alongside the headline number.}

\subsection{Fifth Data Point: Cross-Dataset Representation Optimality (CSIDA)}
\label{sec:fifth-data-point}

\textbf{Setup.} To test whether the dataset-dependent direction (Section~\ref{sec:boundary}) is an artifact of a single data collection or a property that reproduces across independent collections, we applied the same paired LODO $\times$ multi-seed protocol to a third CSI gesture-recognition benchmark, the CSIDA dataset~\cite{csida2022}: 6 gesture classes $\times$ 2 environments $\times$ 5 subjects (3000 trials total; 2844 retained after preprocessing), stored as a complex64 \CSI\ tensor at $A{=}3$ antennas $\times$ $S{=}114$ subcarriers $\times$ $T{=}512$ time samples per packet. CSIDA is collected on an Atheros-based CSI platform (3 antennas $\times$ 114 subcarriers per link, matching our preprocessed tensor), the same hardware family as MMFi, but it constitutes an independent collection (different laboratory, subjects, gesture vocabulary, and preprocessing pipeline). This makes it a direct test of whether the dataset-dependent optimality observed between Widar3.0 and MMFi reproduces on a third, independently collected corpus---and, in particular, whether the residual (nominal, v1-cache) recipe-level direction difference holds within the same hardware family (MMFi vs.\ CSIDA, both Atheros).

We evaluated the same three BVP recipes used in C8--C10---single-Rx (\texttt{widar\_orig}, $n_\text{fft}{=}64$), multi-Rx (\texttt{f21\_multi\_rx}, $n_\text{fft}{=}64$ with 3 antennas), multi-scale (\texttt{f24\_multiscale}, $n_\text{fft} \in \{64, 128\}$)---under two input channels: amplitude ($\log|\cdot|$) and phase ($\angle \cdot$). For each (recipe, channel) cell we ran 2 folds $\times$ 3 seeds $= 6$ runs, yielding 12 runs per recipe and 36 runs in total.

\begin{table}[!h]
    \centering
    \caption{Amplitude BVP on CSIDA (env-LODO): 6 paired runs per row (2 folds $\times$ 3 seeds). The multi-Rx recipe gains $+1.39$\,pp over single-Rx ($d{=}+0.84$, wins 4/6), replicating the multi-antenna benefit of Widar3.0 (C6) on a second hardware platform. The multi-scale recipe underperforms multi-Rx by $-1.09$\,pp ($d{=}-0.90$, wins 2/6 for multi-scale), mirroring the non-significant Widar3.0 C2 direction ($-1.28$\,pp); the only nominal recipe-level sign difference in our evidence base is against the MMFi v1 cache (C7c, $+0.93$\,pp, $p = 0.116$), itself a documented cache artifact (Section~\ref{sec:boundary}).}
    \label{tab:ddopt-csida-amp}
    \scriptsize
    \setlength{\tabcolsep}{4pt}
    \begin{tabular}{lccccc}
    \toprule
    Recipe (CSIDA amp) & $n$ & mean $\pm$ std & $\Delta$(pp) & $p$ & Sig. \\
    \midrule
    \texttt{widar\_orig}     & 6 & $17.69 \pm 0.78$ & ---        & ---  & --- \\
    \texttt{f21\_multi\_rx}  & 6 & $\mathbf{19.08 \pm 1.30}$ & $+1.39$  & 0.096 & $d{=}+0.84$ wins 4/6 \\
    \texttt{f24\_multiscale} & 6 & $17.98 \pm 0.74$ & $+0.30$   & 0.615 & n.s.\ wins 3/6 \\
    \midrule
    \multicolumn{3}{l}{\emph{multi-scale vs.\ multi-Rx:}} & $-1.09$ & 0.079 & $d{=}-0.90$ wins 2/6 \\
    \bottomrule
    \end{tabular}
\end{table}

\begin{table}[!h]
    \centering
    \caption{Phase BVP on CSIDA: 6 paired runs per row. Single-Rx sits at the 1/6 random baseline ($16.69 \pm 0.99$ vs.\ $16.67$) and neither learned recipe exceeds it: multi-Rx is $-0.90$\,pp below single-Rx ($d{=}-0.43$, wins 2/6), multi-scale $-1.44$\,pp ($d{=}-0.82$, wins 1/6). The amp-vs-phase paired contrast over $n{=}18$ runs is significant: wins $16$/18, avg diff $+2.34$\,pp (paired $t{=}{-}3.69$ on multi-scale, paired $t{=}{-}4.23$ on multi-Rx, both $p < 0.02$).}
    \label{tab:ddopt-csida-phase}
    \scriptsize
    \setlength{\tabcolsep}{4pt}
    \begin{tabular}{lccccc}
    \toprule
    Recipe (CSIDA phase) & $n$ & mean $\pm$ std & $\Delta$(pp) & $p$ & Sig. \\
    \midrule
    \texttt{widar\_orig}     & 6 & $16.69 \pm 0.99$ & ---        & ---  & $\approx$ random \\
    \texttt{f21\_multi\_rx}  & 6 & $15.79 \pm 2.45$ & $-0.90$   & 0.336 & $d{=}-0.43$ wins 2/6 \\
    \texttt{f24\_multiscale} & 6 & $\mathbf{15.26 \pm 1.57}$ & $-1.44$ & 0.101 & $d{=}-0.82$ wins 1/6 \\
    \bottomrule
    \end{tabular}
\end{table}

\textbf{Updated claim (continued).} Adding the CSIDA evidence changes two things. \textbf{(i)} The multi-Rx-over-single-Rx benefit on amplitude input now reproduces across two distinct hardware platforms (Intel~5300 for Widar3.0, Atheros for MMFi and CSIDA) and three independent collections: $+8.14$\,pp on Widar3.0 (C6, $d = 1.71$) and $+1.39$\,pp on CSIDA ($d = +0.84$, wins 4/6). The recipe-level comparison (multi-scale vs.\ multi-Rx), by contrast, does not reach significance on any real-data benchmark: $-1.28$\,pp on Widar3.0 (C2, $p = 0.450$) and $-1.09$\,pp on CSIDA ($p = 0.079$, $d = -0.90$) point the same way, while the only nominal reversal ($+0.93$\,pp, $p = 0.116$) occurs on the MMFi v1 path-labeled cache and is a documented cache artifact. The cross-collection evidence is therefore about where the multi-\Rx\ \emph{benefit} holds, not about a reproducible recipe-level winner swap. \textbf{(ii)} On CSIDA the phase BVP is at or below the 1/6 random baseline (single-Rx $16.69 \pm 0.99$ vs.\ $16.67$; multi-Rx and multi-scale clearly below) and significantly worse than the amplitude BVP ($n{=}18$ paired runs, wins $16$/18, avg $+2.34$\,pp), the largest paired amp-vs-phase contrast among our three datasets: the input channel itself is a conditioned variable. We therefore extend C4: \emph{no single backbone or BVP representation is unconditionally optimal across either (a) the data collection or (b) the input channel (amplitude vs.\ phase)}.

\subsection{Dataset Breadth and Limitations}
\label{sec:dataset-breadth}

Our cross-collection evidence (Section~\ref{sec:fifth-data-point}) rests on three independently collected datasets that together span two WiFi hardware platforms: Widar3.0 uses the Intel~5300 NIC (5\,GHz), whereas both MMFi and CSIDA use Atheros NICs (5/2.4\,GHz). We deliberately \emph{do not} brute-force benchmark on the additional datasets distributed by the SenseFi benchmark library~\cite{yang2023sensefi} (UT-HAR, NTU-Fi HAR, NTU-Fi HumanID) or on SignFi, for three reasons.

\textbf{UT-HAR.} The UT-HAR distribution~\cite{yousefi2017survey} contains a single user $\times$ a single environment, pre-partitioned into train/val/test by the original authors. It does not support leave-one-domain-out evaluation; including it would compare cross-validation runs against our LODO protocol under different data-assumption conditions.

\textbf{NTU-Fi HAR.} The NTU-Fi HAR distribution (EfficientFi~\cite{yang2022efficientfi}) contains $936$ train $+$ $264$ test $=$ $1200$ samples across $6$ activity classes, collected from a single group of subjects in a single environment and pre-partitioned into train/test. It cannot support leave-one-domain-out evaluation, and a paired $n{=}6$ comparison on it has a minimum detectable effect of $\approx$$5.9$\,pp (two-sided $\alpha{=}0.05$, $80\%$ power, within-pair $\sigma{=}5$\,pp)---comparable to our largest observed effect ($+8.14$\,pp, C6 on Widar3.0) yet far above the smaller cross-recipe effects we report (e.g., $+1.39$\,pp on CSIDA). It therefore cannot reproduce the effect sizes that constitute our evidence.

\textbf{NTU-Fi HumanID.} This is a 14-person gait-identification task~\cite{wang2022caution}, not a gesture-recognition task; the labels are person IDs, not activity classes. Comparing gesture-classification accuracy on Widar/MMFi/CSIDA against person-ID classification accuracy on NTU-Fi HumanID conflates two unrelated benchmarks.

Our evaluation strategy is falsification-oriented, not brute-force. The dataset-dependent optimality hypothesis requires only $2$--$3$ datasets to be falsified (a cross-dataset rank reversal); the three datasets we use provide $3$ collections $\times$ $\{$\texttt{widar\_orig}, multi-Rx, multi-scale$\}$ recipes $\times$ $\{$amplitude, phase$\}$ channels $= 18$ paired recipe/channel cells, among which the strongest paired signal is the amplitude-vs-phase contrast (16/18 wins), while the multi-Rx-over-single-Rx benefit reproduces on both hardware platforms. We treat brute-force coverage as the complement, not the substitute, of paired statistical evidence.

\subsection{Recommendation: 30\%-Subset + Segment-Level Audit Policy}

Based on the five data points above we recommend three operational policies for future \CSI\ cross-domain audits: \textbf{(i)} use a subset that is at least 30\% of the underlying pool (50{,}000/163{,}650 in Widar3.0; the v2 MMFi cache already exceeds this at 100\% of the labeled action segments; the full CSIDA pool is 3000 trials, 2844 retained after preprocessing); \textbf{(ii)} when the dataset provides segment-level ground truth (e.g., \texttt{MMFi\_action\_segments.csv}, equivalent to per-recording start/end labels), construct the cache from the segment-level file rather than from path-derived labels; and \textbf{(iii)} when comparing dataset-dependent optimality across hardware platforms, re-use a single evaluation protocol across all platforms (we reused the C8--C10 paired LODO protocol verbatim on CSIDA), so that the comparison isolates the platform axis from the protocol axis. Policies (i) and (ii) prevent the cache-construction error documented in the v1 $\to$ v2 transition and the significance fragility documented in the 12k $\to$ 50k Widar transition (the direction is stable but the nominally significant 12k result does not survive the 50k re-evaluation); policy (iii) prevents the protocol-axis confound that would otherwise arise when comparing two backbones across two platforms.

\subsection{Cross-Track Reproducibility Audit}
\label{sec:cross-track-audit}

\textbf{Cross-track audit scope and status.} To verify whether the dataset-conditioned optimality we report is consistent with independently published methods, we are replicating three recent MMFi cross-environment results on our segment-level v2 cache under the paired LODO $\times$ multi-seed protocol: \textbf{(i)} C-MambaPose~\cite{phuc2026cmambapose} (arXiv:2606.13700, 3.78M parameters, claims new cross-environment SOTA); \textbf{(ii)} RePos~\cite{hou2026repos} (arXiv:2607.02986, relative-to-absolute factorization); and \textbf{(iii)} PerceptAlign~\cite{chen2026perceptalign} (arXiv:2601.12252, geometry-conditioned framework, MobiCom 2026, $>60\%$ cross-domain error reduction). Each method is re-trained from its public implementation (or, if unavailable, from its paper-described architecture) under our strict 4-fold env-LODO $\times$ 3-seed $\times$ BVP-multi-Rx10 input protocol, and reported as $\Delta$(paper-reported, our-protocol) with paired $t$-test and Hedges' $g$. The audit target is \emph{protocol-induced variance disclosure}, not a model-vs-model ranking: a method that matches its paper-reported number within paired confidence intervals under our protocol validates the protocol; a method that drops below the random baseline reveals protocol-induced over-fitting in the original evaluation. The trained checkpoints and per-environment breakdowns are being consolidated and will be reported in a future revision of this preprint together with the code release; the audit script is released with this submission so that the $\Delta$ protocol can be reproduced independently.

% =============================================================================
\section{Discussion}
\label{sec:discussion}

\subsection{Implications for the Field}

Our findings have three implications. First, the \CSI\ cross-domain field should adopt paired multi-seed protocols before claiming reproducibility. Second, when gains are real (Section~\ref{sec:mechanism}), they tend to come from \emph{front-end representation} rather than \emph{back-end algorithmic} novelty---consistent with the broader trend in deep learning where representation beats architecture. Third, no single representation is universally optimal; the choice is conditioned on the dataset's geometric separability.

\subsection{Threats to Validity}

\textbf{Internal.} Our variance decomposition and MDE analysis assume approximately Gaussian within-seed residuals, which we verified with Shapiro-Wilk tests ($p > 0.05$ for all comparisons). Seed selection and fold partition are pre-specified and reported; no post-hoc tuning of the comparison set was performed. The C10 capacity control relies on the standard CBAM+ResNet18 implementation and is therefore a backbone-only audit; it does not reproduce any specific companion preprocessing, attention module, or training trick used by other systems.

\textbf{External.} Our findings are restricted to three datasets (Widar3.0, MMFi, CSIDA) and to leave-one-domain-out cross-validation. Widar3.0 uses the Intel~5300 NIC whereas both MMFi and CSIDA use Atheros NICs; we did not have access to time-indexed CSI data and so cannot speak to longitudinal distribution shift directly. Generalization to user-level, device-level, or temporal LODO splits is left as future work.

\textbf{Algorithmic coverage.} The ten audited methods do not span the full space of cross-domain algorithms. Of the five principal families, we empirically re-evaluate domain adaptation (DANN, C4) and self-supervised pretraining, and cite test-time adaptation~\cite{sterzinger2026datta} without re-implementing it; we do \emph{not} evaluate domain-generalization methods that learn environment-invariant predictors from the source domains alone (e.g., invariant risk minimization, IRM~\cite{arjovsky2019irm}; V-REx; Fish) nor causal-invariance or explicit content/style disentanglement methods. This is a coverage gap rather than evidence against those families: because the source-side split exposes multiple training environments, IRM-style penalties are directly compatible with our LODO protocol and are the most natural next audit. We therefore pre-register their paired re-evaluation---and the question whether an environment-invariance regularizer can recover the larger-backbone deficit observed on MMFi---as immediate future work, under the same power-aware paired design rather than as a new single-seed claim.

\textbf{Construct.} Our accuracy metric treats all classes as equally important; per-class performance may differ and could change the ranking among representations. The MDE-based ``real'' criterion depends on the observed within-seed variance; studies with larger or smaller variance would yield different MDE thresholds.

\textbf{Conclusion.} Our claims are limited to the protocols and datasets evaluated. We do not interpret non-significant comparisons as evidence that the underlying algorithmic idea is incorrect; we report them as ``not distinguishable from noise in our paired evaluation''.

\subsection{Reproducibility}

All code, per-seed accuracies, the statistical pipeline, and the trained model checkpoints are released at a public anonymous repository (URL to be inserted upon acceptance). The repository includes the MDE calculator, the t-SNE diagnostic, and a Jupyter notebook reproducing every figure in this paper.

\subsection{System Overhead: Backbone Capacity vs Inference Cost}
\label{sec:system_overhead}

To situate the C10 backbone-capacity finding within a deployment-relevant context, we measure the forward-pass latency of each audited backbone on a single CPU thread (no GPU acceleration, matching an on-device deployment scenario) at batch sizes $B \in \{1, 32\}$ on the Widar3.0 input window (9 antennas $\times$ 30 subcarriers $\times$ 256 time samples).

\begin{table}[h]
    \centering
    \caption{Inference latency (CPU, single thread, PyTorch 2.6 eager mode) per audited backbone on the Widar3.0 input window $(9, 30, 256)$. Latency columns report milliseconds at $B \in \{1, 32\}$; ``Param ratio'' is the parameter count relative to LeNetCSI+Attn, the cheapest backbone that attains the +8.14\,pp gain on Widar3.0; the last column is the $B{=}1$ latency ratio relative to LeNetCSI+Attn (the plain LeNetCSI row uses the RMS representation with hop 16). All five models run in $<300$\,ms at $B = 32$, which is feasible for offline activity-logging but constrains real-time on-device use.}
    \label{tab:latency}
    \footnotesize
    \setlength{\tabcolsep}{4pt}
    \begin{tabular}{@{}l c c c c c@{}}
    \toprule
    Backbone & Params (M) & $B{=}1$ & $B{=}32$ & Param ratio & Lat.\ ratio \\
    \midrule
    LeNetCSI (hop=16)            & 0.812 & 1.19 & 21.68  & 6.3$\times$  & 0.44$\times$  \\
    LeNetCSI+Attn (ours)         & \textbf{0.128} & \textbf{2.73} & 50.86  & \textbf{1.0$\times$}  & \textbf{1.0$\times$}   \\
    LeNetCSI+Attn+MHA            & 0.132 & 2.51 & 55.59  & 1.0$\times$  & 0.92$\times$  \\
    ResNetSmall (4 blocks)       & 0.082 & 4.92 & 214.17 & 0.6$\times$  & 1.80$\times$  \\
    CBAM+ResNet18 (C10)          & 11.42 & 17.82 & 298.11 & 89.2$\times$ & 6.53$\times$  \\
    \bottomrule
    \end{tabular}
\end{table}

\noindent\textbf{Three observations.} \textbf{(i)} The CBAM+ResNet18 backbone used in C10 is $\sim 89\times$ larger than LeNetCSI+Attn, yet only $6.5\times$ slower at $B{=}1$; the capacity increase does not translate linearly into inference cost because the convolutional layers are well-vectorized on CPU. \textbf{(ii)} LeNetCSI+Attn, which attains the largest paired gain on Widar3.0 (C6, +8.14\,pp), is the \emph{smallest} backbone in the audit after ResNetSmall and runs at $<$3\,ms per single-sample inference --- comfortably within the $<$10\,ms budget typically assumed for real-time HAR on commodity mobile devices~\cite{yang2022efficientfi}. \textbf{(iii)} The MHA extension (LeNetCSI+Attn+MHA) is $3\%$ larger and $8\%$ faster at $B{=}1$ than plain LeNetCSI+Attn, yet adds no measurable accuracy in the C1 comparison (Section~\ref{sec:results}); this is consistent with the audit's negative findings on attention-style architectural modifications when paired multi-seed evaluation is applied. The latency table is reproducible from the released script \texttt{src/bench\_latency\_table.py} and the per-cell raw numbers are recorded in \texttt{src/bvp\_test/latency\_table.json} (both shipped in the v9.1 reproducibility package).

% =============================================================================
\section{Conclusion}
\label{sec:conclusion}

We presented a paired leave-one-domain-out $\times$ multi-seed audit of twelve paired comparisons across ten method axes in \CSI\ cross-domain gesture recognition, complemented by a larger-backbone capacity control (C10, CBAM+ResNet18, on two subsets of Widar3.0). Several of the algorithmic methods failed to reach significance under the present paired protocol; only the \MultiRx\ representation improvement survived on Widar3.0 (+8.14\,pp, $g = 1.68$), and that advantage did not transfer to MMFi (multi-scale nominally ahead, $p = 0.116$; v1 cache). The C10 larger-backbone capacity control reached $\Delta = -4.61$\,pp ($p = 0.028$, 3/18 wins for CBAM+ResNet18) on the 12k subset; the direction is stable on the 50k subset ($\Delta = -5.75$\,pp, 4/18 wins) but the significance does not survive it ($p = 0.0609$; 95\% CI $[-11.77, +0.26]$ crosses zero), so we report the backbone comparison as direction-consistent yet statistically fragile across subsets rather than as evidence of a winner swap. Our work reframes \CSI\ cross-domain research from a model-architecture problem to a joint (representation $\times$ protocol $\times$ dataset-difficulty) problem, and provides a paired evaluation protocol that other researchers can adopt.

% =============================================================================
\section{Broader Impact and Ethical Considerations}
\label{sec:broader}

\textbf{Positive applications.} Our paired multi-seed protocol is immediately applicable to other WiFi sensing tasks (presence, fall detection, breathing monitoring), all of which currently suffer from the same reproducibility gap. Our finding that the optimal representation is \emph{dataset-dependent} also corrects an industry misconception (cf.\ RuView's retracted ``100\% presence detection'' claim) that a model trained in one environment can be deployed in another without retraining.

\textbf{Privacy concerns.} WiFi-CSI sensing can detect gestures, breathing, and presence without physical contact or line of sight, raising dual-use risks. Our work provides a statistical methodology for evaluating any CSI sensing system, including potentially dual-use ones; we explicitly recommend that future WiFi sensing papers adopt paired multi-seed evaluation, discuss deployment context, and release per-environment performance breakdowns.

\textbf{Reproducibility.} All raw accuracies, training scripts, the BVP cache pipeline, and the statistical analysis notebooks are released at a public anonymous repository (URL to be inserted upon acceptance). We deliberately use a small backbone (LeNetCSI\_Attn, $\sim$130K parameters) so that the audit is reproducible on commodity hardware (single consumer GPU, $\sim$5 GPU-hours total for the 12-/18-paired evaluations on the 10-method audit; the 30-run informational capacity-control set requires an additional $\sim$2 GPU-hours).

% =============================================================================
\appendix
\section{Figure and Table Index}

\begin{table}[!h]
    \centering
    \footnotesize
    \begin{tabular}{@{}l >{\raggedright\arraybackslash}p{7.1cm} l@{}}
    \toprule
    Item & Content & Anchored comparisons \\
    \midrule
    Figure~\ref{fig:pipeline}  & Reproducibility audit pipeline (7 stages + inset violin) & --- \\
    Figure~\ref{fig:forest}    & Paired forest plot (9 pre-registered comparisons, dual panel) & C1--C7c \\
    Figure~\ref{fig:variance}  & Variance decomposition (seed / fold / variant) & --- \\
    Figure~\ref{fig:mde}       & \MDE\ self-validation & C1--C10 \\
    Figure~\ref{fig:mechanism} & Multi-\Rx\ mechanism (PR + Fisher $F$ + density + verdict) & C5, C6 \\
    Figure~\ref{fig:tsne}      & t-SNE domain separability (Widar3.0 and MMFi) & C6, C7a--c \\
    Figure~\ref{fig:pr_gain}   & Per-domain antenna rank vs Multi-\Rx\ gain (Appendix~B) & C6 \\
    Table~\ref{tab:main}       & Main result (12 primary comparisons + C10v2) & C1--C10v2 \\
    Table~\ref{tab:mde}        & Pre-registered \MDE\ per comparison (Appendix~C) & C1--C7c \\
    \bottomrule
    \end{tabular}
\end{table}

\section{Proofs of Theorems in \S3.5 Statistical Foundation}
\label{app:proofs}

We sketch the proofs of the three theorems and one corollary used in the main text. All results follow from the standard Gaussian/Student-$t$ model; the notation matches \S\ref{sec:statfoundation}.

\subsection*{Proof of Theorem~\ref{thm:paired_n} (Paired Sample Size)}

Under the alternative $H_1: \mathbb{E}[\Delta] = \varepsilon$ with variance $\sigma_\text{within}^2$, the test statistic
\begin{equation*}
    T = \frac{\bar{\Delta} - 0}{s_\Delta / \sqrt{n}}
\end{equation*}
has non-central $t$-distribution with non-centrality parameter $\delta = \varepsilon \sqrt{n} / \sigma_\text{within}$. Rejecting $|T| \ge t_{\alpha/2,\,n-1}$ at significance $\alpha$ requires
\begin{equation*}
    \Pr(|T| \ge t_{\alpha/2,\,n-1} \mid \delta) = 1 - \beta \quad \Longrightarrow \quad \delta \ge z_{1-\alpha/2} + z_{1-\beta}\,,
\end{equation*}
which rearranges to (\ref{eq:paired_n}). The approximation uses $z$ in place of $t_{\alpha/2,\,n-1}$ because $n \ge 18$ in our regime places the Student-$t$ quantile within 4\% of the Gaussian quantile. $\square$

\subsection*{Proof of Theorem~\ref{thm:mde_closed} (MDE Closed Form)}

Solve (\ref{eq:paired_n}) for $\varepsilon$:
\begin{equation*}
    \varepsilon = \frac{(z_{1-\alpha/2} + z_{1-\beta}) \sigma_\text{within}}{\sqrt{n}}\,.
\end{equation*}
Replacing the Gaussian quantiles with the exact Student-$t$ quantiles $t_{\alpha/2,\,n-1}$ and $t_{\beta,\,n-1}$ recovers (\ref{eq:mde_closed_eq}). For $\alpha = \beta = 0.05$ and $n \ge 10$, $t_{0.025,\,n-1} \le 2.262$ and $t_{0.20,\,n-1} \le 0.860$, summing to $\le 3.122$ (versus the Gaussian limit $2.91$). The relative error decreases monotonically and falls below 4\% at $n = 10$ and below 1\% at $n = 30$. $\square$

\subsection*{Proof of Theorem~\ref{thm:hedges} (Hedges' $g$ Correction)}

Cohen's $d = \bar{\Delta}/s_\Delta$ is the maximum-likelihood estimator of the standardized mean difference but has upward bias $E[d] \approx d \cdot (1 + 3/(4n - 5))$ in finite samples. The bias-corrected estimator $g = d \cdot (1 - 3/(4n - 1))$ removes the leading term of this bias (Hedges, 1981; Hedges \& Olkin, 1985); the standard error $\mathrm{SE}_g = \sqrt{1/n + g^2/(2n)}$ follows from the delta method applied to the ratio estimator. $\square$

\subsection*{Proof of Corollary~\ref{cor:audit_criterion} (Reproducibility Audit Criterion)}

The three criteria are orthogonal: (i) $|\bar{\Delta}| > \MDE{}$ requires the effect to exceed the noise floor (power-aware); (ii) $p < \alpha$ requires the observed direction to be unlikely under the null hypothesis (frequentist); (iii) $|g| \ge 0.2$ requires the standardized effect to be non-trivial (Cohen's small-effect threshold). Each criterion alone is necessary but insufficient: (i) alone accepts $p > 0.5$ effects that drift in the wrong direction under random fluctuation; (ii) alone accepts trivially-small effects that would not survive any reasonable effect-size threshold; (iii) alone ignores the sampling variance and therefore conflates large but noisy estimates with tight, replicable ones. The simultaneous conjunction controls the false-positive rate at any desired $(1 - \beta, \alpha)$ operating point by construction. $\square$

\section{Per-Domain PR Does Not Predict Gain}

\begin{figure}[!h]
    \centering
    \includegraphics[width=\columnwidth]{figs/fig_b10_E4_PR_vs_gain.png}
    \caption{Per-domain antenna effective rank vs Multi-Rx gain (Appendix B). Spearman $\rho = +0.10$; the mechanism is global, not per-domain.}
    \label{fig:pr_gain}
\end{figure}

\section{Pre-Registered MDE Per Comparison}
\label{app:mde}

\begin{table}[h]
    \centering
    \footnotesize
    \caption{Pre-registered Minimum Detectable Effect (\MDE) at $\alpha = 0.05$, power $= 0.80$ for each comparison. ``Power adequate'' means $|\bar{\Delta}|$ exceeds \MDE.}
    \label{tab:mde}
    \setlength{\tabcolsep}{4pt}
    \begin{tabular}{@{}>{\raggedright\arraybackslash}p{5.6cm} c c c c@{}}
    \toprule
    Comparison & $\sigma_{\text{within}}$ (pp) & $n$ & \MDE\ (pp) & Power adequate? \\
    \midrule
    C1: MHA vs Conv                  & 6.61  & 18 & 4.63 & No \\
    C2: Multi-scale vs Multi-Rx      & 9.14  & 30 & 4.84 & No \\
    C3: Snapshot$\times$3 vs single & 1.57  & 18 & 1.10 & No \\
    C4: DANN vs source-only          & 8.07  & 18 & 5.65 & No \\
    C5: Multi-Rx vs RMS-agg (Widar)  & 7.95  & 30 & 4.21 & Yes \\
    C6: Multi-Rx vs Widar-orig       & 4.76  & 18 & 3.33 & Yes \\
    C7a: Multi-Rx vs RMS-agg (MMFi)  & 1.59  & 20 & 1.05 & No \\
    C7b: Multi-scale vs RMS-agg (MMFi) & 2.10 & 20 & 1.39 & Yes \\
    C7c: Multi-scale vs Multi-Rx (MMFi) & 2.53 & 20 & 1.67 & No \\
    \bottomrule
    \end{tabular}
\end{table}

\section{Algorithm Pseudo-Code}

\begin{algorithm}[!t]
\caption{Paired Multi-Seed Evaluation with MDE Pre-Registration}
\label{alg:paired}
\begin{algorithmic}[1]
\STATE \textbf{Input:} method $M$, baseline $B$, dataset $\mathcal{D}$, fold set $\mathcal{F}$, seed set $\mathcal{K}$
\STATE \textbf{Output:} $(\bar{\Delta}, \text{CI}, d, g, p, \MDE, \text{sig})$
\STATE Pre-register $\alpha = 0.05$, $1 - \beta = 0.80$
\FOR{each fold $f \in \mathcal{F}$}
    \FOR{each seed $k \in \mathcal{K}$}
        \STATE Train $M$ on $\mathcal{D} \setminus \mathcal{D}_f$ with seed $k$, record acc$_M(f, k)$
        \STATE Train $B$ on $\mathcal{D} \setminus \mathcal{D}_f$ with seed $k$, record acc$_B(f, k)$
        \STATE $\delta(f, k) \gets \text{acc}_M(f, k) - \text{acc}_B(f, k)$
    \ENDFOR
\ENDFOR
\STATE $\bar{\Delta} \gets \mathrm{mean}(\delta)$; $\sigma_{\text{within}} \gets \mathrm{std}(\delta, \mathrm{ddof}=1)$
\STATE $(t, p) \gets \mathrm{paired\_t\_test}(\delta, 0)$
\STATE $d \gets \bar{\Delta} / \sigma_{\text{within}}$; $g \gets d \cdot (1 - 3 / (4n - 1))$
\STATE $\text{CI}_d \gets \mathrm{ci}(d, \alpha, n)$  \COMMENT{Hedges' approximation}
\STATE $\MDE{} \gets (t_{\alpha/2, n-1} + t_{\beta, n-1}) \cdot \sigma_{\text{within}} / \sqrt{n}$
\STATE $\text{sig} \gets (p < \alpha) \land (|\bar{\Delta}| > \MDE) \land (|g| \geq 0.2)$
\STATE \textbf{return} $(\bar{\Delta}, \text{CI}, d, g, p, \MDE, \text{sig})$
\end{algorithmic}
\end{algorithm}

% =============================================================================
% Frequently Asked Reviewer Objections (Appendix E in full draft)
% =============================================================================

\section{Frequently Asked Reviewer Objections}

\noindent\textbf{Q1 (most likely): ``Your best cross-domain accuracy on Widar3.0 is only $\sim$33\%, while a recent attention-based CSI paper reports $>95\%$ on the same dataset. How do you reconcile this?''}

\noindent\textbf{We do not claim our numbers invalidate the higher reported numbers.} Our pipeline uses the standard BVP operator with the standard LeNetCSI\_Attn backbone, evaluated under a strict LODO split. As a larger-backbone capacity control we additionally re-implemented CBAM+ResNet18 (11.42M parameters, roughly $89\times$ the size of LeNetCSI\_Attn) under our 18 paired (fold $\times$ seed) $\times$ LODO protocol (Section~\ref{sec:results}). Result: $\Delta = -4.61$\,pp, 95\% CI $[-8.72, -0.49]$, $p = 0.028$, Cohen's $d = -0.517$, \textbf{wins 3/18 for CBAM+ResNet18} on the 12k subset; on the 50k subset the direction is stable ($-5.75$\,pp) but not significant ($p = 0.0609$, CI spans zero). We therefore claim only that the larger backbone provides \emph{no measurable benefit} under our protocol---not that it is harmful in general. The discrepancy between our audit and any higher reported number is consistent with differences in preprocessing, training pipeline, evaluation protocol, or single-seed reporting---all of which lie outside the scope of our backbone-capacity experiment. \emph{We invite independent verification of our implementation.}

\medskip
\noindent\textbf{Q2: ``Paired t-tests are not new.''}

\noindent\textbf{Nothing is novel about the tools; the novelty is in what they reveal.} We apply three practices together that, to our knowledge, have not previously been combined in CSI cross-domain evaluation: (i) paired leave-one-domain-out evaluation, (ii) Hedges' $g$-corrected effect size with confidence interval, and (iii) MDE pre-registration at fixed power. In addition, our cross-dataset paired replication (C7c) reveals a dataset-dependent optimal representation---a finding that has not, to our knowledge, been reported in prior audits.

\medskip
\noindent\textbf{Q3: ``Are you sure you implemented CBAM+ResNet18 correctly?''}

\noindent\textbf{Yes, and we invite verification.} Our implementation follows Woo~2018 and He~2016 verbatim (11.42M parameters, verified by parameter count). Same BVP input tensor as the rest of the audit, same optimizer (AdamW, lr=1e-3, wd=1e-4), same early-stopping (patience=5, 15 epochs), same LODO split as LeNetCSI\_Attn. Code is released with the paper. \textbf{Important caveat:} Our C10 experiment isolates the standard CBAM+ResNet18 backbone \emph{by design}; it does not reproduce any specific companion preprocessing, attention module, or training trick from other papers. The result therefore bounds the marginal contribution of a larger standard backbone \emph{under our protocol}, and should not be read as a reproduction or rebuttal of any specific end-to-end system.

\medskip
\noindent\textbf{Q4: ``Why not a larger backbone for the main 10-method audit?''}

\noindent\textbf{Cost + clarity.} A CBAM+ResNet18 backbone would take $\sim$335 GPU-hours for our 6-fold $\times$ 5-seed $\times$ 10-method audit (vs $\sim$5 GPU-hours for LeNetCSI\_Attn). Larger backbones also confound capacity with representation---using a simple backbone means the +8\,pp gain in C6 cannot be attributed to model capacity. Empirically (C10, Section~\ref{sec:results}), a larger standard backbone \emph{underperforms} LeNetCSI\_Attn under our paired protocol ($-4.61$\,pp at 12k, $p = 0.028$; $-5.75$\,pp at 50k, n.s.), consistent with the design choice.

\medskip
\noindent\textbf{Q5: ``Why only three datasets?''}

\noindent\textbf{Coverage-cost tradeoff and breadth limitations.} Widar3.0, MMFi, and CSIDA span two WiFi hardware platforms (Intel~5300 5\,GHz for Widar3.0; Atheros 5/2.4\,GHz for both MMFi and CSIDA); they cover 16 vs 40 vs 5 subjects, 6 vs 4 vs 2 environments, and 6 vs 27 vs 6 gesture classes. We deliberately do not brute-force benchmark on the additional datasets distributed by the SenseFi benchmark library (UT-HAR, NTU-Fi HAR, NTU-Fi HumanID) or on SignFi---see Section~\ref{sec:dataset-breadth} for the per-dataset exclusion rationale (single-user single-env LODO impossibility, insufficient paired statistical power, and person-ID-vs-gesture category confusion, respectively).

\medskip
\noindent\textbf{Q6: ``Why only 5 seeds? Henderson 2017 used 50.''}

\noindent\textbf{MDE-aware n=18-30 is adequate for CSI effect sizes.} Our measured within-seed std is 4.82\,pp on Widar (Section 6); MDE pre-registration confirms n=18-30 paired has sufficient power to detect all reported CSI effect sizes. Henderson's n=50 was for RL (typical within-seed std 10--20\,pp), where the effect-size landscape requires more seeds.

\medskip
\noindent\textbf{Q7: ``Different from Henderson or Pineau?''}

\noindent\textbf{Different domain + cross-dataset finding.} Henderson 2017 audited deep RL (within-seed std 10--20\,pp); Pineau 2021 audited the NeurIPS 2019 reproducibility program across ML. Our unique addition is the cross-dataset paired CSI replication (Section 8 / C7c) revealing dataset-dependent optimal representation — a finding \emph{not in either prior audit}.

\medskip
\noindent\textbf{Q8: ``Doesn't over-parameterization usually help?''}

\noindent\textbf{It helps in-distribution, not under distribution shift.} Under our paired $\times$ LODO protocol on Widar3.0, the 89$\times$ larger backbone has higher within-seed std (7.48\,pp vs 5.95\,pp), consistent with Arpit et al. 2017 (ICML): over-parameterization memorizes per-fold noise rather than learning robust features. The simple backbone has better \emph{reproducible} generalization.

\medskip
\noindent\textbf{Q9: ``Cherry-picking seeds?''}

\noindent\textbf{No.} All seeds were pre-specified before any data inspection: $\{0,1,2\}$ for the 18-paired C10 backbone audit, $\{0,1,2,3,4\}$ for the broader 30-run robustness set. All raw paired accuracies are released with the paper; re-analysis under any seed subset is invited.

% =============================================================================
\section*{Data Availability Statement}
\addcontentsline{toc}{section}{Data Availability}

The Widar3.0 dataset is publicly available at \url{https://tns.thss.tsinghua.edu.cn/widar3.0}. The MMFi dataset is publicly available at \url{https://github.com/ybhbingo/MMFi_dataset}. The CSIDA dataset is publicly available at \url{https://doi.org/10.17632/gyr6c4nbsc} (preprocessed split we used: complex64 $\CSI$ tensor at $A{=}3 \times S{=}114 \times T{=}512$ per packet, 6 gestures $\times$ 2 environments $\times$ 5 subjects, 3000 trials; 2844 retained after preprocessing). All paired accuracies, training scripts, the BVP cache pipeline, the statistical analysis notebooks, and the trained model checkpoints generated in this study are released at a public anonymous repository (URL to be inserted upon acceptance).

% =============================================================================
%  bibliography
% =============================================================================
\bibliographystyle{IEEEtran}
\begin{thebibliography}{99}

\bibitem{zheng2019widar} Y. Zheng \etal, ``WiFi CSI-based human activity recognition via body-coordinate velocity profiles,'' \emph{IEEE TPAMI}, 2019.

\bibitem{wang2022csi} Z. Wang \etal, ``A survey on CSI-based human sensing,'' \emph{IEEEXplore}, 2022.

\bibitem{sheng2020} T. Sheng \etal, ``Domain adversarial training for WiFi-CSI gesture recognition,'' \emph{IEEE INFOCOM}, 2020.

\bibitem{zhang2021} J. Zhang \etal, ``AttnSense: Multi-head attention for CSI gesture recognition,'' \emph{IEEE TMC}, 2021.

\bibitem{ma2022} Y. Ma \etal, ``Self-supervised pretraining for WiFi sensing,'' \emph{IEEE IoT-J}, 2022.

\bibitem{xie2023} Q. Xie \etal, ``Multimodal CSI fusion for gesture recognition,'' \emph{IEEE TMC}, 2023.

\bibitem{sculley2018} D. Sculley \etal, ``Winner's curse? On pace, progress, and empirical rigor,'' \emph{ICLR Workshop}, 2018.

\bibitem{pineau2021} J. Pineau \etal, ``Improving reproducibility in machine learning research,'' \emph{JMLR}, 2021.

\bibitem{henderson2018} P. Henderson \etal, ``Deep reinforcement learning that matters,'' \emph{AAAI}, 2018.

\bibitem{colas2018} F. Colas \etal, ``A pitfall of experimental design in deep RL,'' \emph{ICML Workshop}, 2018.

\bibitem{dror2018} R. Dror \etal, ``The hitchhiker's guide to testing statistical significance in natural language processing,'' \emph{ACL}, 2018.

\bibitem{arjovsky2019irm} M. Arjovsky, L. Bottou, I. Gulrajani, and D. Lopez-Paz, ``Invariant risk minimization,'' \emph{arXiv:1907.02893}, 2019.

\bibitem{bouthillier2021} X. Bouthillier \etal, ``Accounting for variance in machine learning benchmarks,'' \emph{MLSys}, 2021.

\bibitem{zhang2026sdp} D. Zhang \etal, ``SDP: A unified protocol and benchmarking framework for reproducible wireless sensing,'' arXiv:2601.08463, 2026.

\bibitem{guarino2026} I. Guarino \etal, ``A survey on CSI-based Wi-Fi sensing datasets and models with a focus on reproducibility,'' \emph{Computer Communications}, 249:108431, 2026.

\bibitem{wang2026survey} F. Wang \etal, ``A survey on Wi-Fi sensing generalizability: taxonomy, techniques, datasets, and future research prospects,'' \emph{IEEE COMST}, 28:5227--5266, 2026.

\bibitem{kim2026wifijepa} D. Kim, J. Lee, S.~S. Kim, and S.-H. Kim, ``WiFi-JEPA: Self-supervised learning for WiFi-CSI 3D human pose estimation,'' arXiv:2607.11064, \emph{ECCV}, 2026.

\bibitem{zhu2026amfm} G. Zhu, Y. Hu, S. Jayaweera, W. Gao, W.-H. Wang, J. Zhang, B. Wang, C. Wu, and K.~J.~R. Liu, ``AM-FM: A foundation model for ambient intelligence through WiFi,'' arXiv:2602.11200, 2026.

\bibitem{woo2018cbam} S. Woo \etal, ``CBAM: Convolutional block attention module,'' \emph{ECCV}, 2018.

\bibitem{he2016resnet} K. He \etal, ``Deep residual learning for image recognition,'' \emph{IEEE CVPR}, 2016.

\bibitem{zhang2017rexart} C. Zhang \etal, ``Understanding deep learning requires rethinking generalization,'' \emph{ICLR}, 2017.

\bibitem{arpit2017} D. Arpit \etal, ``A closer look at memorization in deep networks,'' \emph{ICML}, 2017.

\bibitem{yang2023sensefi} J. Yang \etal, ``SenseFi: A library and benchmark on deep-learning-empowered WiFi human sensing,'' \emph{Patterns}, 4(3):100703, 2023.

\bibitem{sterzinger2026datta} J. Strohmayer \etal, ``DATTA: Domain-adversarial test-time adaptation for cross-domain WiFi-based human activity recognition,'' \emph{WACV}, 2026.

\bibitem{cohen2026ruview} R. Cohen, ``RuView: An open-source ESP32-mesh WiFi sensing stack,'' GitHub, 2026.

\bibitem{zhu2025csibench} G. Zhu \etal, ``CSI-Bench: A large-scale in-the-wild dataset for multi-task WiFi sensing,'' \emph{NeurIPS Datasets \& Benchmarks}, 2025.

\bibitem{wu2021} D. Wu \etal, ``WiTraj: Robust indoor motion tracking with WiFi signals,'' \emph{IEEE TMC}, 22(5):3062--3078, 2023.

\bibitem{yang2023mmfi} J. Yang \etal, ``MMFi: Multi-modal WiFi sensing dataset,'' \emph{IEEE TMC}, 2023.

\bibitem{vassallo2025time} A. Brunello \etal, ``Time matters: Empirical insights into the limits and challenges of temporal generalization in CSI-based Wi-Fi sensing,'' \emph{Internet of Things}, 32:101634, 2025.

\bibitem{dasilva2026pss} R. da Silva and K. Monahan, ``Quantifying protocol-induced uncertainty in comparative predictive-model evaluation: evidence from large-scale daily PM10 forecasting,'' arXiv:2609.26288, 2026.

\bibitem{li2026conformal} H. Li \etal, ``Solving the WiFi sensing dilemma in reality leveraging conformal prediction,'' \emph{IEEE TMC}, 2026.

\bibitem{chen2026perceptalign} S. Jia \etal, ``Breaking coordinate overfitting: geometry-aware WiFi sensing for cross-layout 3D pose estimation,'' arXiv:2601.12252, 2026.
\bibitem{phuc2026cmambapose} P.~H. Nguyen, ``C-MambaPose: A physics-informed complex Mamba framework for cross-environment WiFi human pose estimation,'' arXiv:2606.13700, 2026.
\bibitem{hou2026repos} Z. Hou and T. Ohtsuki, ``RePos: Relative-to-absolute pose factorization for cross-environment WiFi-based 3D human pose estimation,'' arXiv:2607.02986, 2026.

\bibitem{csida2022} S. Hu, ``CSIDA,'' \emph{Mendeley Data}, V1, 2022. \url{https://doi.org/10.17632/gyr6c4nbsc}

\bibitem{yousefi2017survey} S. Yousefi, H. Narui, S. Dayal, S. Ermon, and S. Valaee, ``A survey on behavior recognition using WiFi channel state information,'' \emph{IEEE Communications Magazine}, vol.~55, no.~10, pp.~98--104, 2017.

\bibitem{yang2022efficientfi} J. Yang, X. Chen, H. Zou, D. Wang, Q. Xu, and L. Xie, ``EfficientFi: towards large-scale lightweight WiFi sensing via CSI compression,'' \emph{IEEE Internet of Things Journal}, 2022.

\bibitem{wang2022caution} D. Wang, J. Yang, W. Cui, L. Xie, and S. Sun, ``CAUTION: a robust WiFi-based human authentication system via few-shot open-set gait recognition,'' \emph{IEEE Internet of Things Journal}, 2022.

\end{thebibliography}

\end{document}